% !TeX root = ../../Book3.tex
% 纲要标题：## 附录 H　代数曲线初步
% 状态：新增正文初稿，尚未正式审核。

\chapter{代数曲线初步}
\label{b3:app:H}
\BookChapterNumber{b3:app:H}

\begin{chapterabstract}
	曲线上的局部阶数把有理函数的零点和极点组织成除子，关联线丛则把指定极点界转为截面空间。本附录建立这些构造与有限曲线态射的联系，计算截面维数和一次上同调，并说明典范线丛如何通过对偶配对控制它们。射影直线提供全部计算都可写出的模型；亏格一曲线则由次数一的除子类得到群律，进一步给出三次方程、同构不变量和同源的例子。
\end{chapterabstract}

\begin{learninggoals}
	\begin{itemize}
		\item 用 DVR 的局部阶数计算除子与有限态射的纤维重数，区别总次数和不同纤维点的个数。
		\item 在除子、线丛与有理截面之间转换，核对极点界和关联层的符号。
		\item 用亏格、次数和典范除子计算线丛截面维数，判断线性系何时无基点或给出嵌入。
		\item 在射影直线的两个仿射图上计算上同调和对偶配对，说明全局截面为何不保持满射。
		\item 从亏格一曲线构造三次模型，通过除子类群验证割线、切线所给群运算的结合性。
		\item 计算短 Weierstrass 方程的判别式与 \(j\) 不变量，并在同源算例中区分次数、点核与不可分性。
	\end{itemize}
\end{learninggoals}

给定射影直线上的两个不同点 \(P,Q\)，可以找到有理函数，使它只在 \(P\) 有一个单零点、在 \(Q\) 有一个单极点。但把同样的要求放在一条光滑三次曲线上，通常就没有解。零点与极点各自的总次数相等，仍不足以保证函数存在。这一差别要求我们同时记录局部阶数、整体的除子类，以及给定极点条件下究竟有多少个线性无关的函数。

\ProseParagraphBreak
除子把有限组局部整数放在同一个对象中；线丛说明这些条件怎样在开集上黏合。截面空间的维数随除子次数增加而增长，但在低次数处会受到曲线本身的影响。亏格与典范线丛给出描述这种影响的两个量。一次上同调所记录的黏合障碍通过一个对偶配对化为正则微分，截面维数公式因而同时包含函数与微分。

\ProseParagraphBreak
本附录使用第6章的正规化有限性、第10章的 DVR、第14章的微分模、第20章的除子与可逆模，以及第23章的光滑局部模型。第24章的仿射模层对应、Proj 和 Čech 比较使这些代数构造成为曲线上的层与上同调；其中一般仿射上同调消失仍沿用该章注明出处的证明概要。第二卷第4章的 PID 模结构定理用于两张仿射图上的有限维数证明。这里从代数闭域上的光滑射影整曲线开始；前三节允许任意特征，最后一节的短 Weierstrass 方程假定特征不为 \(2,3\)。

\ProseParagraphBreak
一般曲线的 Serre 对偶采用有明确技术步骤和出处的证明概要，射影直线的配对则直接算出。其余除子计算、Euler 特征递推和椭圆曲线群律均在相应条件下展开。进一步的曲线理论可参阅\cite[Chapter IV, Sections 1--4]{Hartshorne1977}。

% !TeX root = ../../Book3.tex
% 纲要标题：### H.1 曲线、除子与亏格
% 状态：新增正文初稿，尚未正式审核。

\section{曲线、除子与亏格}
\label{b3:sec:H01}

在仿射直线上，多项式只有零点；补入无穷远点以后，次数又成为那一点的极点阶数。例如坐标函数 \(t\) 在零点消失一次，在无穷远点有一阶极点。一般曲线未必有一个坐标足以表示全部函数，但每个光滑点仍有自己的统一化元。把这些局部阶数放在一起，就得到除子；规定函数所允许的极点，则得到一个线丛的截面空间。

\ProseParagraphBreak
本节在代数闭域 \(k\) 上工作，不限制特征。代数闭性使每个闭点的剩余域都等于 \(k\)，因而除子的次数可直接相加。若底域不代数闭，次数须按 \([\kappa(P):k]\) 加权，且一个次数一的除子类未必由有理点表示；后面关于亏格一曲线的点与除子类的对应依赖这里的基域条件。

\subsection{代数曲线与函数域}
\label{b3:subsec:H01:01}

\begin{definition}\label{b3:def:H-smooth-projective-curve}
	设 \(k\) 为代数闭域。一个一维、有限型、分离的整 \(k\) 概形称为本附录所讨论的\term[daishu-quxian]{代数曲线}[algebraic curve]。若它在 \(k\) 上光滑，并且存在闭嵌入 \(C\hookrightarrow\mathbb P_k^N\)，则称为光滑射影整曲线。其一般点记为 \(\eta\)，\(k(C)=\mathcal O_{C,\eta}\) 是它的\term[quxian-hanshuyu]{函数域}[function field]。
\end{definition}

任取非空仿射开集 \(U=\operatorname{Spec}A\)，便有 \(k(C)=\operatorname{Frac}(A)\)。这里 \(A\) 是有限型一维整环，故函数域在 \(k\) 上的超越次数为一。函数域中的元素是定义在某个非空开集上的正则函数，两个表示在共同定义域内相等就表示同一个有理函数。

\ProseParagraphBreak
代数闭域是完美域，\cref{b3:thm:perfect-field-regular-smooth}因而使光滑曲线在闭点 \(P\) 的局部环成为一维正则局部环。其极大理想由一个元素生成，\cref{b3:thm:dvr-characterizations}使 \(\mathcal O_{C,P}\) 成为 DVR。记相应归一化赋值为 \(\operatorname{ord}_P:k(C)^*\rightarrow\mathbb Z\)。它对函数给出零点或极点的阶数；统一化元的选择只改变单位因子。
\symidx{\(k(C)\)：整曲线 \(C\) 的函数域}
\symidx{\(\operatorname{ord}_P\)：光滑曲线闭点 \(P\) 处的阶数}

\begin{lemma}\label{b3:lem:H-rational-map-extension}
	设 \(k\) 为代数闭域，\(X\) 为正规整曲线，\(Y\subseteq\mathbb P_k^N\) 为射影闭子概形。每个从 \(X\) 到 \(Y\) 的有理 \(k\) 映射都唯一延拓为 \(k\) 态射 \(X\rightarrow Y\)。
\end{lemma}
\begin{proof}
	在一般点把映射写成齐次坐标 \([a_0:\cdots:a_N]\)，其中 \(a_i\in k(X)\)，且不全为零。给定闭点 \(P\)，其局部环为 DVR。取统一化元 \(t\)，令 \(m\) 为非零 \(a_i\) 的阶数最小值。将全部坐标乘以 \(t^{-m}\)，则每个坐标都在 \(\mathcal O_{X,P}\) 中，且至少一个为单位。缩小 \(P\) 的开邻域，使这些坐标都正则，选定的单位坐标处处不为零，就得到到射影空间的态射。\(Y\) 的齐次方程在函数域中成立，所以在这些邻域上仍成立，像落在 \(Y\)。

	两份局部延拓在一般点相等。因为 \(Y\) 分离，它们的等化子是闭子概形；其定义理想在整的源上泛点为零，因而为零。于是两态射相等，可唯一黏合。一般点原来就在定义域中，所有闭点又已覆盖，得到全体 \(X\) 上的延拓。
\end{proof}

这个坐标论证也见\cite[Chapter I, Proposition 6.8, pp. 43--44]{Hartshorne1977}。源必须有能把全部分母统一比较的局部参数；在奇异点，局部环未必是 DVR，向射影空间的有理映射就不能照搬这一消分母过程。

\begin{proposition}\label{b3:prop:H-curve-morphism}
	设 \(k\) 为代数闭域，\(C,D\) 为光滑射影整曲线，\(\varphi:C\rightarrow D\) 为非恒定 \(k\) 态射。则 \(\varphi\) 是有限、平坦、满射的态射。令
	\[
		d=[k(C):\varphi^*k(D)].
	\]
	对任一闭点 \(Q\in D\)，取其统一化元 \(u\)，则有
	\[
		\sum_{P\in\varphi^{-1}(Q)}\operatorname{ord}_P(\varphi^*u)=d.
	\]
	这个整数 \(d\) 称为态射的次数，包含函数域扩张的可分次数与不可分次数。若 \(d=1\)，则 \(\varphi\) 为同构。
\end{proposition}
\begin{proof}
	非恒定性使像稠密：一维整概形的真闭子集只有有限个闭点，整源的像若落在其中就只能是一个点。因此有函数域嵌入 \(k(D)\hookrightarrow k(C)\)。两域都在 \(k\) 上有限生成且超越次数一，所以此扩张代数且有限。

	取 \(D\) 的非空仿射开集 \(V=\operatorname{Spec}A\)，并令 \(B\) 为 \(A\) 在 \(k(C)\) 中的整闭包。\cref{b3:thm:normalization-finite}适用于有限型 \(k\) 整环 \(A\) 及这个有限域扩张，故 \(B\) 是有限 \(A\) 代数。有限域扩张中的元素清分母后，其某个非零 \(A\) 倍数整于 \(A\)，所以 \(\operatorname{Frac}(B)=k(C)\)。整扩张保持维数，故它是一维正规整环，\(\operatorname{Spec}B\) 的闭点局部环都是 DVR。一般点给出到 \(C\) 的有理映射，由上一引理延拓为 \(\psi:\operatorname{Spec}B\rightarrow C\)。复合 \(\varphi\psi\) 与自然映射 \(\operatorname{Spec}B\rightarrow D\) 在一般点相同，由分离性相同，故 \(\psi\) 的像落在 \(\varphi^{-1}(V)\)。

	反向，对 \(P\in\varphi^{-1}(V)\)，每个 \(b\in B\) 都整于 \(A\subseteq\mathcal O_{C,P}\)，故整于 \(\mathcal O_{C,P}\)。该局部环整闭，因此 \(b\in\mathcal O_{C,P}\)。这些有理函数在每一点正则，可黏合成 \(\varphi^{-1}(V)\) 上的正则函数，给出态射 \(\gamma:\varphi^{-1}(V)\rightarrow\operatorname{Spec}B\)。两复合都在一般点为恒等，再由分离性和源的整性，均为恒等。所以
	\[
		\varphi^{-1}(V)\cong\operatorname{Spec}B.
	\]
	这证明 \(\varphi\) 有限；整扩张的 lying-over 定理保证它在每张 \(V\) 上满射，故全局满射。

	在闭点 \(Q\) 处局部化，\(B_Q=B\otimes_A\mathcal O_{D,Q}\) 是 DVR \(\mathcal O_{D,Q}\) 上秩为 \(d\) 的有限生成无挠模，所以自由秩为 \(d\)。这一局部结论可直接用 Nakayama 引理检验：取模统一化元的一组基并提升为最小生成元；若存在非零关系，把系数中共同的最低幂约去，再模掉统一化元，就与这些剩余类线性无关矛盾。一般点处的局部映射是域扩张，也平坦。因此 \(\varphi\) 平坦，且
	\[
		\dim_k B_Q/uB_Q=d.
	\]
	这个有限维 \(k\) 代数分解为各闭点 \(P\mid Q\) 的局部因子的积。每个因子是 \(\mathcal O_{C,P}/(\varphi^*u)\)，其 \(k\) 维数由\cref{b3:prop:dvr-order-ideals}等于 \(\operatorname{ord}_P(\varphi^*u)\)。维数相加便得公式。最后，若 \(d=1\)，函数域相同；\(A\) 已整闭，所以 \(B=A\)，上述仿射识别说明 \(\varphi\) 是同构。
\end{proof}

有限态射与纤维重数公式也见\cite[Chapter II, Propositions 6.8--6.9, pp. 137--138]{Hartshorne1977}。例如特征 \(p>0\) 时，\(\mathbb P^1\rightarrow\mathbb P^1\)、\(t\mapsto t^p\) 的次数是 \(p\)。每个闭点纤维只有一个点，因为 \(t^p-a=(t-b)^p\)，其中 \(b^p=a\)；其重数却为 \(p\)。因此一般不能用不同纤维点的个数代替次数。

\subsection{除子与线丛}
\label{b3:subsec:H01:02}

\begin{definition}\label{b3:def:H-curve-divisor}
	设 \(k\) 为代数闭域，\(C\) 为光滑射影整曲线。一个除子是闭点的有限整数线性组合 \(D=\sum_P n_PP\)，次数为 \(\deg D=\sum_Pn_P\)。各系数非负时称为有效除子，记 \(D\ge0\)。对 \(f\in k(C)^*\)，定义
	\[
		\operatorname{div}(f)=\sum_P\operatorname{ord}_P(f)P.
	\]
	这样的除子称为主除子；若 \(D-E=\operatorname{div}(f)\)，则记 \(D\sim E\)，称 \(D,E\) 线性等价。
\end{definition}

主除子的支撑确实有限。取有限个仿射开集覆盖 \(C\)，在每张图上把 \(f\) 写成分式 \(a/b\)。非零的 \(a,b\) 在一维 Noether 整环中的零点只有有限个：包含它们的素理想都极小于相应主理想，而且都为闭点。有限张图的这些有限集合包含全部非零阶数点。赋值的加法公式使主除子组成除子群的子群，这与\secref{b3:sec:2003}的仿射定义相容。

\begin{corollary}\label{b3:cor:H-principal-degree}
	设 \(k\) 为代数闭域，\(C\) 为光滑射影整曲线。则每个主除子次数为零，并且 \(H^0(C,\mathcal O_C)=k\)。
\end{corollary}
\begin{proof}
	若 \(f\in k(C)^*\) 非恒定，它定义到 \(\mathbb P_k^1\) 的有理映射，由\cref{b3:lem:H-rational-map-extension}延拓为态射。零点纤维与无穷远纤维的总重数，由\cref{b3:prop:H-curve-morphism}都等于 \([k(C):k(f)]\)。前者组成 \(f\) 的零除子，后者组成极除子，故 \(\deg\operatorname{div}(f)=0\)。常函数的主除子为零，同样成立。

	全局正则函数没有极点。若它非恒定，相应映射 \(C\rightarrow\mathbb P^1\) 却必须满射，因而在无穷远点有非空纤维，与没有极点矛盾。因此全局函数都在 \(k\) 中；反向常函数当然全局正则。
\end{proof}

在 \(\mathbb P_k^1\) 上以 \(t\) 为仿射坐标，记无穷远点为 \(\infty\)。对不同的 \(a,b\in k\)，有
\[
	\operatorname{div}\!\left(\frac{(t-a)^2}{t-b}\right)=2[a]-[b]-[\infty].
\]
无穷远处令 \(s=t^{-1}\)，原函数为 \(s^{-1}(1-as)^2/(1-bs)\)，所以该处极点阶数为一。若只保留仿射直线，就漏掉这个负项；次数为零是射影曲线上的全局结论。

\begin{definition}\label{b3:def:H-divisor-sheaf}
	设 \(k\) 为代数闭域，\(C\) 为光滑射影整曲线，\(D=\sum n_PP\) 为除子。对非空开集 \(U\subseteq C\)，定义
	\[
		\mathcal O_C(D)(U)=\{0\}\cup\{\,f\in k(C)^*: \operatorname{ord}_P(f)+n_P\ge0\ \text{对所有闭点}\;P\in U\,\}.
	\]
	记 \(L(D)=H^0(C,\mathcal O_C(D))\)、\(\ell(D)=\dim_kL(D)\)。\(L(D)\) 是满足指定零点及极点界的函数空间。
\end{definition}
\symidx{\(\mathcal O_C(D)\)：除子 \(D\) 的关联线丛}
\symidx{\(L(D)\)、\(\ell(D)\)：除子截面空间及其维数}

这里采用几何的符号约定：正的 \(n_P\) 允许极点，而负的 \(n_P\) 要求零点。它与\secref{b3:sec:2005}中的分式理想 \(I_D\) 相差一个负号；那里的 \(I_D\) 对应这里的 \(\mathcal O_C(-D)\)。

\begin{proposition}\label{b3:prop:H-divisor-line-bundle}
	设 \(k\) 为代数闭域，\(C\) 为光滑射影整曲线，\(D,E\) 为 \(C\) 上的除子。则 \(\mathcal O_C(D)\) 为可逆层，且
	\[
		\mathcal O_C(D)\otimes\mathcal O_C(E)\cong\mathcal O_C(D+E).
	\]
	映射 \([D]\mapsto[\mathcal O_C(D)]\) 给出除子模线性等价所成的群与线丛同构类群 \(\operatorname{Pic}(C)\) 的同构。特别地，线丛的次数可定义为 \(\deg\mathcal O_C(D)=\deg D\)，与表示无关。
\end{proposition}
\begin{proof}
	在 \(P\) 处取统一化元 \(t\)，则所定义层的茎为 \(t^{-n_P}\mathcal O_{C,P}\)。缩小邻域，删去 \(D\) 的其他支撑点以及 \(t\) 的其他零点、极点，这个式子就给出邻域上的自由基。条件由局部正则性定义，因而满足限制与黏合。两个局部基相乘得到 \(t^{-n_P-m_P}\)，证明张量积公式。

	任取线丛 \(\mathcal L\)。其一般点纤维是一维 \(k(C)\) 空间，选非零有理截面 \(s\)。在有限组平凡化开集上写 \(s=f_ie_i\)，其中 \(e_i\) 是局部基。定义 \(n_P=\operatorname{ord}_P(f_i)\)；更换基只乘以单位，所以与选择的开集无关。各 \(f_i\) 的零极点有限，故 \(D_s=\sum n_PP\) 是除子。映射 \(h\mapsto hs\) 把 \(\mathcal O_C(D_s)\) 同构到 \(\mathcal L\)：在 \(P\) 处，\(t^{-n_P}s\) 是一个局部基。这证明满射性。

	若两线丛 \(\mathcal O_C(D),\mathcal O_C(E)\) 同构，其一般点上的映射是乘以某个 \(a\in k(C)^*\)。在 \(P\) 处比较两边的生成元阶数，得 \(\operatorname{ord}_P(a)-n_P(D)=-n_P(E)\)，即 \(D-E=\operatorname{div}(a)\)。反之，此等式使乘以 \(a\) 给出同构。因此核恰是主除子群。次数的良定义由\cref{b3:cor:H-principal-degree}得到。
\end{proof}

这也说明线丛上的非零正则截面给出有效零除子，且其次数等于线丛的次数；对 \(\mathcal O_C(D)\) 的非零截面 \(f\)，零除子正是 \(\operatorname{div}(f)+D\)。除子与可逆层的比较见\cite[Chapter II, Proposition 6.13, Proposition 6.15 and Corollary 6.16, pp. 144--145]{Hartshorne1977}。

\ProseParagraphBreak
次数是群同态 \(\deg:\operatorname{Pic}(C)\rightarrow\mathbb Z\)。其核记为 \(\operatorname{Pic}^0(C)\)，由零次数线丛的同构类组成，也可以写成零次数除子模主除子所成的群。这里把它作为抽象交换群使用；在亏格一时，\secref{b3:sec:H04}将把这个群结构转移到曲线的点集上。
\symidx{\(\operatorname{Pic}^0(C)\)：曲线 \(C\) 的零次数线丛类群}

\begin{corollary}\label{b3:cor:H-section-degree}
	设 \(k\) 为代数闭域，\(C\) 为光滑射影整曲线，\(\mathcal L\) 为线丛。若 \(\deg\mathcal L<0\)，则 \(H^0(C,\mathcal L)=0\)。若 \(\deg\mathcal L=0\)，则它有非零全局截面当且仅当 \(\mathcal L\cong\mathcal O_C\)，此时截面空间维数为一。
\end{corollary}
\begin{proof}
	非零截面的零除子有效，次数不可能为负。次数零时，这个有效除子只能为零，所以该截面处处生成线丛，给出 \(\mathcal O_C\cong\mathcal L\)。反向使用常数截面，再由\cref{b3:cor:H-principal-degree}计算维数。
\end{proof}

在射影直线上，\(\operatorname{div}(t-a)=[a]-[\infty]\)，所以每个除子都线性等价于 \((\deg D)[\infty]\)。不同次数不能线性等价，故 \(\operatorname{Pic}(\mathbb P^1)\cong\mathbb Z\)。在两张坐标图上，\(\mathcal O(n[\infty])\) 的基可取 \(1,t^n\)，其过渡关系正是\secref{b3:sec:2403}中 \(\mathcal O(n)\) 的过渡关系，因而二者同构。

\subsection{亏格与典范除子}
\label{b3:subsec:H01:03}

截面空间给出可以全局实现的函数，一次上同调则保留局部截面之间无法消去的差。为把它的维数作为曲线的数值，需要先确认维数有限。对曲线，这可以直接化为两个 Laurent 多项式格的比较。

\begin{proposition}\label{b3:prop:H-curve-cohomology-finite}
	设 \(k\) 为代数闭域，\(C\) 为光滑射影整曲线，\(\mathcal L\) 为线丛。则 \(H^0(C,\mathcal L)\)、\(H^1(C,\mathcal L)\) 都为有限维 \(k\) 空间，且 \(H^i(C,\mathcal L)=0\) 对 \(i\ge2\) 成立。
\end{proposition}
\begin{proof}
	选一个非恒定函数 \(t\in k(C)\)，它给出次数为 \(d\) 的有限态射 \(\varphi:C\rightarrow\mathbb P^1\)。两张标准图的逆像 \(U,V\) 都仿射，交集 \(W\) 也仿射。设 \(M=\Gamma(U,\mathcal L)\)、\(N=\Gamma(V,\mathcal L)\)。由\cref{b3:cor:vector-bundle-projective-module}，它们分别是两张仿射图坐标环上的有限生成投射模。因为 \(\varphi\) 有限且 \(\mathcal L\) 可逆，\(M,N\) 分别为 \(k[t]\)、\(k[t^{-1}]\) 上的有限生成无挠模；在函数域上，它们都给出一维 \(k(C)\) 空间，故在 \(k(t)\) 上的秩是 \(d\)。

% 跨卷引用目标：b2:thm:pid-module-structure（4.3.2）。
	第二卷定理4.3.2证明主理想整环上有限生成模的不变因子分解；这里两个多项式环都是 PID，无挠性排除全部循环挠商，因此 \(M,N\) 均自由秩为 \(d\)。在交集上限制截面对应于局部化，这是\cref{b3:thm:affine-quasicoherent-equivalence}的模层对应。选 \(M\) 的基，在 \(W\) 上得到 \(E=k[t,t^{-1}]^d\)。\(N\) 的基在其中由某个 \(T\in\operatorname{GL}_d(k[t,t^{-1}])\) 给出，所以
	\[
		M=k[t]^d,\qquad N=T\,k[t^{-1}]^d\quad\text{均作为}\;E\;\text{的子模}.
	\]
	由\cref{b3:thm:affine-cech-comparison}，上同调是两项复形 \(M\oplus N\rightarrow E\)、\((m,n)\mapsto n-m\) 的上同调。因此 \(H^0\cong M\cap N\)、\(H^1\cong E/(M+N)\)，更高次数为零。

	取非负整数 \(r\)，使 \(T\) 及 \(T^{-1}\) 的所有 Laurent 项指数均不超过 \(r\)。若 \(v\in M\cap N\)，写 \(v=Tb\)、\(b\in k[t^{-1}]^d\)，则 \(v\) 的最高指数不超过 \(r\)；又因 \(v\in M\)，其最低指数非负。于是 \(H^0\) 嵌入由 \(t^je_i\)、\(0\le j\le r\) 张成的有限维空间。对 \(H^1\)，所有 \(j\ge0\) 的 \(t^je_i\) 已属于 \(M\)；若 \(j\le-r\)，则 \(T^{-1}t^je_i\) 的指数都非正，故 \(t^je_i\in N\)。余核只需由 \(-r<j<0\) 的有限多个向量生成，也有限维。
\end{proof}

\begin{definition}\label{b3:def:H-genus-canonical}
	设 \(k\) 为代数闭域，\(C\) 为光滑射影整曲线。定义它的\term[kuige]{亏格}[genus]为
	\[
		g(C)=\dim_kH^1(C,\mathcal O_C).
	\]
	将各仿射图上的微分模黏合，得到微分层 \(\Omega_{C/k}\)，记为 \(\omega_C\)，称为曲线的\term[dianfan-xiancong]{典范线丛}[canonical line bundle]。任取非零有理微分 \(\alpha\in\Omega_{k(C)/k}\)，其作为 \(\omega_C\) 的有理截面的除子称为\term[dianfan-chuzi]{典范除子}[canonical divisor]，记为 \(K_C\)。
\end{definition}
\symidx{\(g(C)\)：光滑射影整曲线的亏格}
\symidx{\(\omega_C\)、\(K_C\)：典范线丛与典范除子}

微分层的黏合由\cref{b3:prop:differential-localization}保证；\cref{b3:prop:standard-smooth-lifting,b3:thm:smooth-local-standard-form}使其在光滑曲线上局部自由秩一。在闭点处，统一化元 \(t\) 的类生成 \(\mathfrak m_P/\mathfrak m_P^2\)，由\cref{b3:prop:tangent-cotangent-rational}，\(\mathrm{d}t\) 模极大理想非零，故生成 \(\omega_{C,P}\)。写 \(\alpha=a\,\mathrm{d}t\)，便有 \(\operatorname{ord}_P(\alpha)=\operatorname{ord}_P(a)\)。更换统一化元时，两组局部基相差单位，阶数不变。更换非零有理微分 \(\alpha\) 为 \(f\alpha\)，典范除子增加 \(\operatorname{div}(f)\)；因此 \(K_C\) 只在具体代表上有选择，线性等价类固定，且 \(\omega_C\cong\mathcal O_C(K_C)\)。

\ProseParagraphBreak
在正特征中不能把任意非恒定函数的微分当作典范除子的代表。例如 \(\mathrm{d}(t^p)=0\)，零微分没有定义有限阶数的除子。上面的定义明确要求非零有理微分；光滑性保证函数域上的微分空间是一维，因而存在这样的选择。

\begin{example}\label{b3:exm:H-projective-line-canonical}
	设 \(k\) 为代数闭域，\(C=\mathbb P_k^1\)。在有限点 \(a\)，\(t-a\) 为统一化元，且 \(\mathrm{d}t=\mathrm{d}(t-a)\)，所以 \(\mathrm{d}t\) 没有零极点。在无穷远处令 \(s=t^{-1}\)，则
	\[
		\mathrm{d}t=-s^{-2}\,\mathrm{d}s.
	\]
	所以 \(\operatorname{div}(\mathrm{d}t)=-2[\infty]\)，并有 \(\omega_{\mathbb P^1}\cong\mathcal O(-2)\)。这个等式在特征二仍成立，因为系数 \(-1\) 仍为单位。由\cref{b3:prop:projective-line-cohomology}，\(g(\mathbb P^1)=h^1(\mathcal O)=0\)。同时 \(h^0(\omega_{\mathbb P^1})=0\)，这是\secref{b3:sec:H03}中一般对偶关系的一个具体情形。
\end{example}

亏格目前定义为 \(h^1(\mathcal O_C)\)，而 \(h^0(\omega_C)\) 数的是正则微分。它们相等并不是定义的一部分；\cref{b3:thm:H-serre-duality}将通过一个非退化配对建立两者的联系。典范除子的次数也由\secref{b3:sec:H02}的截面维数公式确定。

% BEGIN APPENDIX OUTLINE SYNC
% * 状态：初稿已完成；新增范围尚未正式审核。
% #### H.1.1 代数曲线与函数域
% #### H.1.2 除子与线丛
% #### H.1.3 亏格与典范除子
% END APPENDIX OUTLINE SYNC

% !TeX root = ../../Book3.tex
% 纲要标题：### H.2 Riemann–Roch 定理
% 状态：正文、证明与算例初稿；配套习题由附录入口统一编排。

\section{Riemann–Roch 定理}
\label{b3:sec:H02}

在射影直线上，允许有理函数只在无穷远点有至多 \(n\) 阶极点，得到的正是次数不超过 \(n\) 的多项式；当 \(n\ge0\) 时，这个空间有 \(n+1\) 个自由参数。一般曲线上，极点总阶数相同，函数空间的维数却可能不同。亏格描述其中一个固定的修正，而典范除子还记录低次数时未消失的上同调。

\ProseParagraphBreak
本节仍令 \(k\) 为代数闭域，\(C\) 为光滑射影整曲线，\(g=h^1(C,\mathcal O_C)\)，并沿用\cref{b3:def:H-divisor-sheaf}中的 \(L(D)\) 和 \(\ell(D)=h^0(C,\mathcal O_C(D))\)。先从逐点增加极点的短正合列计算 Euler 特征，再用\cref{b3:thm:H-serre-duality}把一次上同调写成微分的空间，得到 Riemann–Roch 公式。\footnote{罗赫（Gustav Roch，1839-1866），德国数学家，与 Riemann 的工作共同形成曲线上除子与函数空间维数的 Riemann–Roch 定理。}后一个结果的证明概要在\secref{b3:sec:H03}中给出；其构造通过射影空间对偶与局部 Koszul 复形完成。

\subsection{线性系与 Riemann–Roch 公式}
\label{b3:subsec:H02:01}

\begin{definition}\label{b3:def:H-linear-system}
	设 \(k\) 为代数闭域，\(C\) 为光滑射影整曲线，\(D\) 为 \(C\) 上的除子。全部与 \(D\) 线性等价的有效除子组成的集合
	\[
		|D|=\{\,E\ge0\mid E\sim D\,\}
	\]
	称为 \(D\) 的\term[wanquan-xianxingxi]{完全线性系}[complete linear system]。非零子空间 \(V\subseteq L(D)\) 中各非零截面的零除子组成一个\term[xianxingxi]{线性系}[linear system]，记为 \(|V|\)。若一个闭点 \(P\) 出现在其中每个有效除子内，就称 \(P\) 为该线性系的\term[jidian]{基点}[base point]；没有基点时称线性系\term[wu-jidian]{无基点}[base-point-free]。
\end{definition}

对 \(0\ne f\in L(D)\)，其作为 \(\mathcal O_C(D)\) 截面的零除子是 \(D+\operatorname{div}(f)\)。两个截面给出同一除子，当且仅当其比的除子为零；由\cref{b3:cor:H-principal-degree}，这个比属于 \(k^\times\)。因此 \(|D|\) 是 \(L(D)\) 中一维子空间组成的射影空间，其维数为 \(\ell(D)-1\)。若 \(\ell(D)=0\)，则 \(|D|=\varnothing\)，约定维数为 \(-1\)。同样，\(|V|\) 的维数为 \(\dim_kV-1\)。

\ProseParagraphBreak
判断基点时须使用线丛的局部基。若 \(D\) 在 \(P\) 处的系数为 \(n\)，统一化元为 \(t\)，则局部基为 \(t^{-n}\)；截面 \(f\) 在 \(P\) 处消失，意思是 \(t^nf\in\mathfrak m_P\)，并不总是指有理函数 \(f\) 的值为零。取 \(V\) 的一组基 \(s_0,\ldots,s_r\)，在没有共同零点的开集 \(U\subseteq C\) 上，局部平凡化给出
\[
	\varphi_V:U\longrightarrow\mathbb P_k^r,\qquad
	P\longmapsto[s_0(P):\cdots:s_r(P)].
\]
更换局部基会把所有坐标同时乘以一个单位，所以映射相容；更换 \(V\) 的基则改变目标的射影坐标。若 \(|V|\) 无基点，这些局部映射在整个 \(C\) 上黏合，且 \(\varphi_V^*\mathcal O_{\mathbb P^r}(1)\cong\mathcal O_C(D)\)。在 \(s_i\ne0\) 的开集上，构造的环同态就是把射影图的坐标送到 \(s_j/s_i\)；这些函数在交叠上满足相同的坐标变换。若完全线性系所定义的态射是闭嵌入，就称 \(\mathcal O_C(D)\) 为\term[hen-chongzu-xiancong]{很充足线丛}[very ample line bundle]，也称 \(D\) 很充足。

\ProseParagraphBreak
维数公式的第一步只需要上同调长正合列。由\cref{b3:prop:H-curve-cohomology-finite}，线丛的零次、一次上同调有限维，更高次为零，因而可以定义
\[
	\chi(\mathcal L)\coloneqq h^0(C,\mathcal L)-h^1(C,\mathcal L).
\]
这里的 \(\chi\) 称为线丛的\term[oulatezheng]{Euler 特征}[Euler characteristic]。

\begin{proposition}\label{b3:prop:H-euler-characteristic-divisor}
	设 \(k\) 为代数闭域，\(C\) 为光滑射影整曲线，\(g=h^1(C,\mathcal O_C)\)。对 \(C\) 上的任意除子 \(D\)，有
	\[
		\chi(\mathcal O_C(D))=\deg D+1-g.
	\]
\end{proposition}
\begin{proof}
	先取一个闭点 \(P\)，并记 \(i_P:\{P\}\hookrightarrow C\)。若 \(D\) 在 \(P\) 处的系数为 \(n\)，则包含 \(\mathcal O_C(D)\subseteq\mathcal O_C(D+P)\) 的茎之商为
	\[
		t^{-n-1}\mathcal O_{C,P}/t^{-n}\mathcal O_{C,P}.
	\]
	它被 \(\mathfrak m_P\) 零化，在 \(k(P)=k\) 上一维；其余闭点处的商为零。因此有短正合列
	\[
		0\longrightarrow\mathcal O_C(D)\longrightarrow\mathcal O_C(D+P)
		\longrightarrow i_{P*}Q_P\longrightarrow0,
	\]
	其中 \(Q_P\) 为一维 \(k\) 空间。选择局部基可把它识别为 \(k\)，但这里不需要指定这种识别。点支撑层 \(i_{P*}Q_P\) 的限制映射满射，故为 flasque 层，零次上同调为 \(Q_P\)，正次上同调为零。\cref{b3:def:sheaf-cohomology}后的长正合列于是成为
	\[
		\begin{aligned}
			0&\longrightarrow H^0(\mathcal O_C(D))\longrightarrow H^0(\mathcal O_C(D+P))
			\longrightarrow Q_P\\
			&\longrightarrow H^1(\mathcal O_C(D))\longrightarrow H^1(\mathcal O_C(D+P))
			\longrightarrow0.
		\end{aligned}
	\]
	对这个有限维正合列取维数的交错和，得到
	\[
		\chi(\mathcal O_C(D+P))=\chi(\mathcal O_C(D))+1.
	\]
	另一方面，\(\deg(D+P)=\deg D+1\)。因而所求等式对 \(D\) 成立，当且仅当对 \(D+P\) 成立。

	在 \(D=0\) 时，\cref{b3:cor:H-principal-degree}给出 \(h^0(\mathcal O_C)=1\)，亏格的定义给出 \(h^1(\mathcal O_C)=g\)，故等式成立。任意除子都可从零除子经过有限次增加或减去闭点得到；上面的等价性同时处理这两种操作，便得全部情形。
\end{proof}

\begin{theorem}[Riemann–Roch]\label{b3:thm:H-riemann-roch}
	设 \(k\) 为代数闭域，\(C\) 为亏格 \(g\) 的光滑射影整曲线，\(K_C\) 为一个非零有理微分的除子。对 \(C\) 上的任意除子 \(D\)，有
	\[
		\ell(D)-\ell(K_C-D)=\deg D+1-g.
	\]
\end{theorem}
\begin{proof}
	\cref{b3:thm:H-serre-duality}用于线丛 \(\mathcal O_C(D)\)，给出
	\[
		H^1(C,\mathcal O_C(D))^\vee
		\cong H^0(C,\omega_C\otimes\mathcal O_C(-D)).
	\]
	由典范除子的定义，右端同构于 \(H^0(C,\mathcal O_C(K_C-D))\)，故一次上同调的维数为 \(\ell(K_C-D)\)。代入\cref{b3:prop:H-euler-characteristic-divisor}即得。
\end{proof}

这就是曲线上的\term[Riemann-Roch-dingli]{Riemann–Roch 定理}[Riemann--Roch theorem]。其 Euler 特征推导与对偶输入可对照\cite[Chapter IV, Theorem 1.3, pp. 295--296]{Hartshorne1977}。左端的第二项并非只由 \(\deg D\) 决定；它在低次数时衡量实际的上同调修正。

\begin{corollary}\label{b3:cor:H-canonical-degree}
	设 \(k\) 为代数闭域，\(C\) 为亏格 \(g\) 的光滑射影整曲线，\(K_C\) 为典范除子。则
	\[
		\ell(K_C)=g,\qquad \deg K_C=2g-2.
	\]
	若 \(g=1\)，则 \(K_C\sim0\)，即 \(\omega_C\cong\mathcal O_C\)。若 \(\deg D>2g-2\)，则
	\[
		h^1(C,\mathcal O_C(D))=0,\qquad \ell(D)=\deg D+1-g.
	\]
\end{corollary}
\begin{proof}
	在 Riemann–Roch 公式中先取 \(D=0\)，由 \(\ell(0)=1\) 得 \(\ell(K_C)=g\)；再取 \(D=K_C\)，得到 \(g-1=\deg K_C+1-g\)，即次数公式。若 \(g=1\)，\(K_C\) 次数为零而有非零截面，\cref{b3:cor:H-section-degree}给出 \(K_C\sim0\)。最后，\(\deg D>2g-2\) 使 \(K_C-D\) 次数为负，同一推论给出 \(\ell(K_C-D)=0\)；对偶与 Riemann–Roch 公式分别给出所列两式。
\end{proof}

\begin{proposition}\label{b3:prop:H-high-degree-linear-systems}
	设 \(k\) 为代数闭域，\(C\) 为亏格 \(g\) 的光滑射影整曲线，\(D\) 为除子。若 \(\deg D\ge2g\)，则 \(|D|\) 无基点；若 \(\deg D\ge2g+1\)，则 \(D\) 很充足。
\end{proposition}
\begin{proof}
	对任一闭点 \(P\)，若 \(\deg D\ge2g\)，则 \(D\) 和 \(D-P\) 的次数都大于 \(2g-2\)。上一推论给出 \(\ell(D)-\ell(D-P)=1\)。而 \(L(D-P)\) 正是 \(L(D)\) 中在 \(P\) 处消失的截面，所以每点都有不消失的截面，得到无基点性。

	若 \(\deg D\ge2g+1\)，对任意两个闭点 \(P,Q\)（允许 \(P=Q\)），同样有
	\[
		\ell(D)-\ell(D-P-Q)=2.
	\]
	这表示求值到长度二子概形 \(P+Q\) 的映射满射。若 \(P\ne Q\)，可分别指定两个纤维中的值，因而分离两点。若 \(P=Q\)，在 \(P\) 处取局部基后，\(\mathcal O_C/\mathfrak m_P^2\) 由 \(1,t\) 张成；可指定常数项和一次项，因而分离切向。射影概形上由截面定义的态射，在同时分离闭点和切向时为闭嵌入，这里采用\cite[Chapter II, Proposition 7.3, pp. 152--153]{Hartshorne1977}的判据。其局部结论是目标局部环到源局部环的满射：切向分离给出极大理想模平方的满射，固有性及有限纤维给出局部环的有限模条件，再用 Nakayama 引理提升到局部环。当前 \(C\) 射影、\(k\) 代数闭，正好满足该判据的假设。因此 \(\varphi_{L(D)}\) 为闭嵌入。
\end{proof}

次数界的曲线表述也见\cite[Chapter IV, Proposition 3.1 and Corollary 3.2, pp. 307--308]{Hartshorne1977}。界 \(2g+1\) 是统一的充分条件；小于这个界的线丛仍可能很充足，须检查具体线性系。

\subsection{射影直线与低亏格曲线的计算}
\label{b3:subsec:H02:02}

先回到 \(C=\mathbb P_k^1\)，在有限图上取坐标 \(t\)，无穷远点记为 \(P_\infty\)。对 \(a\in k\)，有 \(\operatorname{div}(t-a)=P_a-P_\infty\)，因此任意次数 \(n\) 的除子都线性等价于 \(nP_\infty\)。没有有限极点的有理函数是多项式，而无穷远处的极点阶数就是其次数，所以
\[
	L(nP_\infty)=
	\begin{cases}
		\operatorname{span}_k\{1,t,\ldots,t^n\},&n\ge0,\\
		0,&n<0.
	\end{cases}
\]
用 \(u=t^{-1}\) 表示另一张图，\(\mathrm{d}t=-u^{-2}\mathrm{d}u\)，故 \(\operatorname{div}(\mathrm{d}t)=-2P_\infty\)。\cref{b3:prop:projective-line-cohomology}已经由两图 Čech 复形给出 \(g=0\)；上述函数和微分计算逐项核验了 Riemann–Roch 公式。

\begin{proposition}\label{b3:prop:H-genus-zero-line}
	设 \(k\) 为代数闭域，\(C\) 为光滑射影整曲线。则 \(C\) 的亏格为零，当且仅当 \(C\cong\mathbb P_k^1\)。
\end{proposition}
\begin{proof}
	射影直线的亏格已在上面算出。反过来，设 \(g=0\)，取闭点 \(P\)。典范除子的次数为 \(-2\)，故 \(\ell(K_C-P)=0\)，Riemann–Roch 给出 \(\ell(P)=2\)。常数只占一维，所以存在非恒定 \(f\in L(P)\)。它只有 \(P\) 处的一个简单极点；否则它处处正则而为常数。由\cref{b3:prop:H-curve-morphism}，\(f:C\rightarrow\mathbb P_k^1\) 是有限态射，其次数等于无穷远纤维的总重数，因而为一。

	这使两个函数域相同。在目标的任一仿射开集上，有限态射的源环整于目标环，且包含在同一函数域中。目标环正规，故两环相等。逐图黏合得到 \(C\cong\mathbb P_k^1\)。
\end{proof}

\begin{proposition}\label{b3:prop:H-genus-one-sections}
	设 \(k\) 为代数闭域，\(C\) 为亏格一的光滑射影整曲线，\(D\) 为除子。则
	\[
		\ell(D)=
		\begin{cases}
			\deg D,&\deg D>0,\\
			1,&D\sim0,\\
			0,&\deg D=0\;\text{且}\;D\not\sim0,\\
			0,&\deg D<0.
		\end{cases}
	\]
	此外，每个次数一的除子类都恰含一个闭点。
\end{proposition}
\begin{proof}
	由\cref{b3:cor:H-canonical-degree}，\(K_C\sim0\)。若 \(\deg D>0\)，则 \(-D\) 次数为负，Riemann–Roch 给出 \(\ell(D)=\deg D\)。次数为零或负的结论分别来自\cref{b3:cor:H-section-degree}以及 \(H^0(C,\mathcal O_C)=k\)。

	若 \(\deg D=1\)，则 \(\ell(D)=1\)。非零截面的零除子是次数一的有效除子，因 \(k\) 代数闭，它就是一个闭点。所有非零截面彼此只差非零常数，所以该有效除子唯一，得到所述唯一性。
\end{proof}

固定亏格一曲线上的一个点 \(O\) 后，\(L(O)=k\)，但 \(L(2O)\) 与 \(L(3O)\) 的维数分别为二、三。可以依次选取在 \(O\) 处极点阶数恰为二、三的函数；不存在仅有一个简单极点的非恒定函数。这些维数将在\secref{b3:sec:H04}中形成平面三次方程。另一方面，次数为零的线丛即使局部都秩一，也可能没有任何非零全局截面；它是否平凡不能仅靠次数判断。

\begin{example}[亏格二的典范线性系]\label{b3:exm:H-genus-two-canonical}
	设 \(k\) 为代数闭域，\(C\) 为亏格二的光滑射影整曲线。则 \(\deg K_C=2\)、\(\ell(K_C)=2\)。对任一闭点 \(P\)，都有 \(\ell(P)=1\)：若有非恒定函数只在 \(P\) 处有至多一个简单极点，便如\cref{b3:prop:H-genus-zero-line}的证明得到 \(C\cong\mathbb P^1\)，与亏格二矛盾。

	\ProseParagraphBreak
	对 \(D=P\) 用 Riemann–Roch，得到 \(\ell(K_C-P)=\ell(P)=1\)。因此 \(L(K_C-P)\) 在 \(L(K_C)\) 中余维一，典范线性系没有基点，给出态射 \(C\rightarrow\mathbb P^1\)。一个目标点的拉回是典范线性系中的一个除子，次数为二，故这个态射的次数为二。这里没有把它当作闭嵌入：两维截面只给出一维射影目标。
\end{example}


% 新增结构与纲要同步；正文、证明与算例初稿，尚未正式审核。
% #### H.2.1 线性系与 Riemann–Roch 公式
% #### H.2.2 射影直线与低亏格曲线的计算
% * 状态：初稿已完成；新增范围尚未正式审核。

% !TeX root = ../../Book3.tex
% 纲要标题：### H.3 Serre 对偶与上同调计算
% 状态：正文与算例初稿；一般对偶给证明概要，射影直线配对逐项计算。

\section{Serre 对偶与上同调计算}
\label{b3:sec:H03}

在射影直线的两张标准图上，\(t^{-1}\) 代表 \(H^1(\mathcal O(-2))\) 中不能消去的类。把它乘以微分 \(\mathrm{d}t\) 后，取 Laurent 展开中 \(t^{-1}\mathrm{d}t\) 的系数，就得到一个数。一般曲线上的 Serre 对偶把这种计算组织为一次上同调与微分截面之间的完美配对；\secref{b3:sec:H02}所用的修正项 \(\ell(K_C-D)\) 正由此出现。

\ProseParagraphBreak
以下仍设 \(k\) 为代数闭域，\(C\) 为光滑射影整曲线。记 \(\omega_C=\Omega_{C/k}\)，并以 \(V^\vee=\operatorname{Hom}_k(V,k)\) 表示有限维 \(k\) 空间的对偶。一般配对的存在性先用射影嵌入给出证明概要；两图上的配对则直接由截面复形算出。

\subsection{曲线上的对偶配对}
\label{b3:subsec:H03:01}

若 \(\mathcal L\) 是线丛，一个截面 \(s\in H^0(C,\omega_C\otimes\mathcal L^{-1})\) 等价于层同态 \(s:\mathcal L\rightarrow\omega_C\)。它在一次上同调上诱导线性映射。因而只要给出 \(H^1(C,\omega_C)\) 到 \(k\) 的适当迹映射，就可把一次上同调类与 \(s\) 配成一个数。

\begin{theorem}[曲线上的 Serre 对偶]\label{b3:thm:H-serre-duality}
	设 \(k\) 为代数闭域，\(C\) 为光滑射影整曲线。存在迹同构
	\[
		\operatorname{Tr}_C:H^1(C,\omega_C)\longrightarrow k,
		\qquad\omega_C=\Omega_{C/k},
	\]
	使对每个线丛 \(\mathcal L\)，下列配对都是有限维 \(k\) 空间之间的完美配对：
	\[
		\begin{aligned}
			H^1(C,\mathcal L)\times H^0(C,\omega_C\otimes\mathcal L^{-1})
			&\longrightarrow k,\\
			(\alpha,s)&\longmapsto\operatorname{Tr}_C(H^1(s)(\alpha)).
		\end{aligned}
	\]
	换言之，这一配对诱导自然同构
	\[
		H^1(C,\mathcal L)\cong H^0(C,\omega_C\otimes\mathcal L^{-1})^\vee.
	\]
	把 \(\mathcal L\) 换成 \(\omega_C\otimes\mathcal L^{-1}\)，也得到 \(H^0(C,\mathcal L)\) 与 \(H^1(C,\omega_C\otimes\mathcal L^{-1})\) 的完美配对。
\end{theorem}
\begin{proofsketch}
	取闭嵌入 \(i:C\hookrightarrow P=\mathbb P_k^N\)，记 \(c=N-1\)、\(\omega_P=\bigwedge^N\Omega_{P/k}\)。射影空间上的对偶首先给出
	\[
		\operatorname{Ext}^{c}_{\mathcal O_P}(i_*\mathcal L,\omega_P)
		\cong H^1(P,i_*\mathcal L)^\vee
		\cong H^1(C,\mathcal L)^\vee.
	\]
	第一步是\cite[Chapter III, Theorem 7.1, p. 240]{Hartshorne1977}：先对 \(\mathcal O_P(a)\) 用标准仿射覆盖上的单项式配对，再对有限个扭曲层的直和作同样计算；一般凝聚层由这些直和展示，零次的 Hom 配对由正合列延拓，高次 Ext 则由可消去性和导出函子的比较得到。第二步可在 \(P\) 的仿射覆盖上检验：闭嵌入的逆像仍仿射，\(i_*\mathcal L\) 的截面与 \(\mathcal L\) 的截面相同，两份 Čech 复形逐项一致。

	为把左端写回 \(C\) 上，设 \(\mathcal I\) 为 \(C\) 的理想层，取
	\[
		\omega'_C\coloneqq
		\mathcal{E}xt^c_{\mathcal O_P}(i_*\mathcal O_C,\omega_P),
	\]
	把它视作 \(C\) 上的层。在光滑嵌入的局部，\(\mathcal I\) 由长度 \(c\) 的正则序列生成。该序列的 Koszul 复形是 \(\mathcal O_C\) 的局部自由消解，其 Hom 复形在 \(c\) 次以前无上同调。截取内射消解时，这些较低次的正合项可以消去，得到
	\[
		\operatorname{Hom}_C(\mathcal L,\omega'_C)
		\cong\operatorname{Ext}^c_{\mathcal O_P}(i_*\mathcal L,\omega_P).
	\]
	这一 Hom 与 Ext 的变换，以及由它得到的对偶化层构造，分别见\cite[Chapter III, Lemma 7.4 and Proposition 7.5, p. 242]{Hartshorne1977}。其技术内容包括在层范畴中保持内射性的伴随识别，以及低次正合复形的分裂截取。

	同一 Koszul 计算还说明，更换正则序列生成元时，最高外幂带来行列式因子；用余法丛的行列式对偶补偿后，局部同构黏合为
	\[
		\omega'_C\cong\omega_P|_C\otimes\det(\mathcal I/\mathcal I^2)^\vee
		\cong\Omega_{C/k}.
	\]
	最后一个同构来自光滑闭嵌入的余法正合列
	\[
		0\longrightarrow\mathcal I/\mathcal I^2
		\longrightarrow\Omega_{P/k}|_C
		\longrightarrow\Omega_{C/k}\longrightarrow0
	\]
	的最高外幂。这里采用\cite[Chapter III, Theorem 7.11 and Corollary 7.12, pp. 245--246]{Hartshorne1977}的局部完全交计算及典范层识别。

	合成上述同构，得到对 \(\mathcal L\) 自然的
	\[
		\operatorname{Hom}_C(\mathcal L,\omega_C)
		\cong H^1(C,\mathcal L)^\vee.
	\]
	令 \(\mathcal L=\omega_C\)，恒等同态对应一个泛函 \(\operatorname{Tr}_C\)。由自然性，任意 \(s:\mathcal L\rightarrow\omega_C\) 所对应的泛函正是 \(\operatorname{Tr}_C\circ H^1(s)\)，这给出定理中的配对。又因 \(\operatorname{Hom}_C(\omega_C,\omega_C)=H^0(C,\mathcal O_C)=k\)，迹映射为同构。

	这里省略射影空间上一般凝聚层的导出函子比较、光滑嵌入的正则序列构造，以及层 Ext 的截取和黏合细节；上述文献给出了这些步骤的证明。一般局部自由层的对偶形式另见\cite[Chapter III, Corollary 7.7, p. 244; Corollary 7.12 and Remark 7.12.2, p. 246]{Hartshorne1977}。
\end{proofsketch}

这一定理称为\term[Serre-duiou]{Serre 对偶}[Serre duality]。对偶空间中的微分带有 \(\mathcal L^{-1}\)；特别地，当 \(\mathcal L=\mathcal O_C(D)\) 时，它是 \(\omega_C\otimes\mathcal O_C(-D)\cong\mathcal O_C(K_C-D)\)，除子的符号由此确定。

\ProseParagraphBreak
配对还可直接传到 Čech 代表元。对有限仿射开覆盖 \((U_i)\)，取表示 \(\alpha\) 的交叠截面 \(a_{ij}\in\mathcal L(U_i\cap U_j)\)。乘以全局截面 \(s\) 后，\((s a_{ij})\) 是 \(\omega_C\) 的余圈。乘法与限制相容，且把余边送到余边，因此所得上同调类只依赖 \(\alpha,s\)。\cref{b3:thm:affine-cech-comparison}识别这份计算与层上同调；迹随后将它送到 \(k\)。

\ProseParagraphBreak
若 \(u:\mathcal L\rightarrow\mathcal M\) 是线丛同态，则 \(s\circ u\) 与 \(H^1(u)(\alpha)\) 满足
\[
	\langle H^1(u)(\alpha),s\rangle_{\mathcal M}
	=\langle\alpha,s\circ u\rangle_{\mathcal L}.
\]
因而在完美配对的同构下，有交换图
\[
	\begin{tikzcd}[column sep=large]
		H^1(C,\mathcal L)\arrow[r,"H^1(u)"]\arrow[d,"\cong"']
		&H^1(C,\mathcal M)\arrow[d,"\cong"]\\
		H^0(C,\omega_C\otimes\mathcal L^{-1})^\vee
		\arrow[r,"(u^*)^\vee"']
		&H^0(C,\omega_C\otimes\mathcal M^{-1})^\vee .
	\end{tikzcd}
\]
这里 \(u^*\) 表示把到 \(\omega_C\) 的同态与 \(u\) 复合。这个方向说明，上同调映射的对偶是微分截面上的反向映射。

\subsection{线丛上同调与典型计算}
\label{b3:subsec:H03:02}

射影直线上的迹可以明确指定，不必保留为抽象泛函。取
\[
	U_0=\operatorname{Spec}k[t],\qquad
	U_\infty=\operatorname{Spec}k[u],\qquad u=t^{-1}.
\]
对 \(\mathcal O(n)=\mathcal O(nP_\infty)\)，以 \(e_0=1\)、\(e_\infty=t^n\) 为局部基，所以 \(e_\infty=t^ne_0\)。采用“后一张减前一张”的 Čech 微分，截面复形为
\[
	k[t]e_0\oplus k[t^{-1}]e_\infty
	\longrightarrow k[t,t^{-1}]e_0,\qquad
	(ae_0,be_\infty)\longmapsto(t^n b-a)e_0.
\]
两图及交集都仿射，\cref{b3:thm:affine-cech-comparison}保证它计算层上同调。核由 \(a(t)=t^nb(t^{-1})\) 给出；第一边的指数非负，第二边的指数至多为 \(n\)，所以核的基为 \(e_0,te_0,\ldots,t^ne_0\)（\(n\ge0\)），\(n<0\) 时为零。余核则为
\[
	H^1(\mathbb P^1,\mathcal O(n))
	\cong\frac{k[t,t^{-1}]}{k[t]+t^nk[t^{-1}]}\,e_0.
\]
分母消去指数 \(j\ge0\) 和 \(j\le n\)，留下的基为 \([t^je_0]\)、\(n<j<0\)。这同时给出
\[
	h^0(\mathcal O(n))=\max\{n+1,0\},\qquad
	h^1(\mathcal O(n))=\max\{-n-1,0\},
\]
并与\cref{b3:prop:projective-line-cohomology}的计算一致。

\begin{proposition}[射影直线上的显式配对]\label{b3:prop:H-projective-line-duality}
	设 \(k\) 为代数闭域，\(n\in\mathbb Z\)。在上述坐标与局部基下，取迹同构
	\[
		\operatorname{Tr}_{\mathbb P^1}([h(t)\,\mathrm{d}t])
		=\operatorname{coeff}_{t^{-1}}h(t),
	\]
	其中右端表示 Laurent 多项式中 \(t^{-1}\) 的系数。若 \(n\le-2\)，则
	\[
		\begin{aligned}
			H^1(\mathcal O(n))&=\bigoplus_{i=1}^{-n-1}k\cdot[t^{-i}e_0],\\
			H^0(\omega_{\mathbb P^1}\otimes\mathcal O(-n))
			&=\bigoplus_{j=0}^{-n-2}k\,t^j\mathrm{d}t\otimes e_0^{-1},
		\end{aligned}
	\]
	配对为
	\[
		\langle[he_0],q\,\mathrm{d}t\otimes e_0^{-1}\rangle
		=\operatorname{coeff}_{t^{-1}}(hq).
	\]
	特别地，\([t^{-i}e_0]\) 与 \(t^{i-1}\mathrm{d}t\otimes e_0^{-1}\) 互为对偶基。若 \(n\ge-1\)，这两个空间都为零。
\end{proposition}
\begin{proof}
	在第二张图上 \(\mathrm{d}u=-t^{-2}\mathrm{d}t\)，所以正则微分的两图贡献为 \(k[t]\mathrm{d}t\) 与 \(t^{-2}k[t^{-1}]\mathrm{d}t\)。它们都没有 \(t^{-1}\mathrm{d}t\) 项，而商只剩这一项。因此所写迹良定义且为同构。

	对 \(\omega_{\mathbb P^1}\otimes\mathcal O(-n)\)，局部基分别为 \(\mathrm{d}t\otimes e_0^{-1}\) 和 \(\mathrm{d}u\otimes e_\infty^{-1}\)，后者等于前者的 \(-t^{-n-2}\) 倍。全局截面因而写成 \(q(t)\mathrm{d}t\otimes e_0^{-1}\)，其中 \(q\in k[t]\)、\(\deg q\le-n-2\)。若这个上界为负，只有零截面。

	乘积与迹给出所列公式。若把 \(h\) 增加一个 \(a(t)\in k[t]\)，则 \(a q\) 的指数非负，不影响 \(t^{-1}\) 系数；若增加 \(t^n b(t^{-1})\)，乘 \(q\) 后的最高指数至多为 \(n+(-n-2)=-2\)，也不影响该系数。因此配对在所写余核上良定义。最后，对 \(1\le i,r\le-n-1\)，\(t^{-i}t^{r-1}\) 的 \(t^{-1}\) 系数在 \(i=r\) 时为一，否则为零，故为完美配对。\(n\ge-1\) 时两端都为零，结论同样成立。
\end{proof}

在另一方向，当 \(n\ge0\) 时，\(H^0(\mathcal O(n))\) 的基 \(t^ie_0\)（\(0\le i\le n\)）与
\[
	[t^{-i-1}\mathrm{d}t\otimes e_0^{-1}]
	\in H^1(\omega_{\mathbb P^1}\otimes\mathcal O(-n))
\]
互为对偶基。这里的相同系数计算解释了为什么零次截面的非负指数，与另一线丛一次上同调的负指数恰好对应。上述迹的标准 Čech 生成元也见\cite[Chapter III, Remark 7.1.1, pp. 240--241]{Hartshorne1977}。

\ProseParagraphBreak
回到一般曲线，若有效除子 \(Z\) 的次数为 \(m\)，则限制线丛 \(\mathcal L|_Z\) 的全局截面空间为 \(m\) 维：在每个支撑点，长度为该点系数的局部商环上，线丛仍为秩一自由模。短正合列
\[
	0\longrightarrow\mathcal L(-Z)\longrightarrow\mathcal L
	\longrightarrow\mathcal L|_Z\longrightarrow0
\]
给出求值映射及其障碍：
\[
	\begin{aligned}
		0&\longrightarrow H^0(\mathcal L(-Z))\longrightarrow H^0(\mathcal L)
		\xrightarrow{\mathrm{ev}_Z}H^0(\mathcal L|_Z)\\
		&\xrightarrow{\delta}H^1(\mathcal L(-Z))\longrightarrow H^1(\mathcal L)
		\longrightarrow0.
	\end{aligned}
\]
在 \(Z=2P\) 时，求值指定的是局部展开的常数项和一次项。若 \(\deg\mathcal L-m>2g-2\)，\cref{b3:cor:H-canonical-degree}使 \(H^1(\mathcal L(-Z))=0\)，所以这些局部条件都能同时满足；次数较小时，\(\delta\) 正记录不能满足的条件。Serre 对偶把这一障碍空间识别为微分截面空间的对偶。

\begin{example}[次数为零的两种上同调]\label{b3:exm:H-degree-zero-cohomology}
	设 \(k\) 为代数闭域，\(C\) 为亏格 \(g\ge1\) 的光滑射影整曲线，\(P,Q\) 为不同闭点。线丛 \(\mathcal O_C(P-Q)\) 不平凡：若 \(P-Q\) 为主除子，就存在只有一个简单极点的非恒定函数，由\cref{b3:prop:H-curve-morphism}得到到 \(\mathbb P^1\) 的次数一态射，进而 \(C\cong\mathbb P^1\)，与 \(g\ge1\) 矛盾。由\cref{b3:cor:H-section-degree}，该线丛的零次上同调为零，而 Euler 特征公式给出
	\[
		\begin{aligned}
			\bigl(h^0(\mathcal O_C),h^1(\mathcal O_C)\bigr)&=(1,g),\\
			\bigl(h^0(\mathcal O_C(P-Q)),h^1(\mathcal O_C(P-Q))\bigr)&=(0,g-1).
		\end{aligned}
	\]
	两者次数均为零，上同调却不同。特别地，亏格一曲线上的非平凡零次数线丛没有任何上同调；亏格至少二时，其一次上同调仍非零。
\end{example}

\begin{proposition}\label{b3:prop:H-genus-one-cohomology}
	设 \(k\) 为代数闭域，\(C\) 为亏格一的光滑射影整曲线，\(\mathcal L\) 为次数 \(d\) 的线丛。则
	\[
		\bigl(h^0(C,\mathcal L),h^1(C,\mathcal L)\bigr)=
		\begin{cases}
			(d,0),&d>0,\\
			(1,1),&\mathcal L\cong\mathcal O_C,\\
			(0,0),&d=0\;\text{且}\;\mathcal L\not\cong\mathcal O_C,\\
			(0,-d),&d<0.
		\end{cases}
	\]
\end{proposition}
\begin{proof}
	\cref{b3:cor:H-canonical-degree}给出 \(\omega_C\cong\mathcal O_C\)，所以 Serre 对偶把 \(H^1(C,\mathcal L)\) 识别为 \(H^0(C,\mathcal L^{-1})^\vee\)。任意线丛由\cref{b3:prop:H-divisor-line-bundle}写成 \(\mathcal O_C(D)\)，再对 \(D\) 和 \(-D\) 分别应用\cref{b3:prop:H-genus-one-sections}，便得到所列四种情形。
\end{proof}

这张维数表中，正次数的消失与负次数的非消失互为对偶；零次数还必须区分平凡与非平凡线丛。一般亏格下也有同样的计算方式：先辨认 \(K_C-D\) 的有效截面，再由对偶确定一次上同调，最后用 Euler 特征恢复零次维数。次数界能使某些项直接消失，低次数时则需要具体的除子类信息。


% 新增结构与纲要同步；正文与算例初稿，尚未正式审核。
% #### H.3.1 曲线上的对偶配对
% #### H.3.2 线丛上同调与典型计算
% * 状态：初稿已完成；新增范围尚未正式审核。

% !TeX root = ../../Book3.tex
% 纲要标题：### H.4 椭圆曲线的几何
% 内容：亏格一曲线的方程、由除子类定义的群律及同源的基本计算。

\section{椭圆曲线的几何}
\label{b3:sec:H04}

亏格为一时，正次数除子的截面数恰好等于次数。因此，允许函数在一个指定点有二阶、三阶极点，便能选出两个函数；它们之间最早出现的关系给出三次方程。另一方面，一次除子类的唯一点代表又把除子类的加法变成曲线上点的运算。方程与群律由此来自同一组 Riemann--Roch 计算。

\ProseParagraphBreak
本节仍以代数闭域 \(k\) 上的光滑射影整曲线为对象。由除子类定义点的群运算不限制特征；短 Weierstrass 方程、以下坐标公式及同源计算均假设 \(\operatorname{char}k\ne2,3\)。方程构造与群结构可对照\cite[Chapter IV, Proposition 4.6, Proposition 4.8, pp. 319--322]{Hartshorne1977}。

\subsection{Weierstrass 方程与光滑平面三次曲线}
\label{b3:subsec:H04:01}

\begin{definition}\label{b3:def:H-elliptic-curve}
	设 \(k\) 为代数闭域，\(E\) 为光滑射影整曲线，\(O\in E(k)\) 为一个指定点。若 \(g(E)=1\)，则称 \((E,O)\) 为一条\term[tuoyuan-quxian]{椭圆曲线}[elliptic curve]。指定点 \(O\) 称为原点。椭圆曲线的同构要求是把指定原点送到指定原点的 \(k\) 曲线同构。
\end{definition}

这里的亏格是 \(h^1(E,\mathcal O_E)\)，并不是方程中某个变量的次数。原点也不是额外的一组坐标；它将确定函数可以出现极点的位置，以及后面的群单位元。

\begin{theorem}[Weierstrass 模型]\label{b3:thm:H-weierstrass-model}
	设 \(k\) 为代数闭域，\(\operatorname{char}k\ne2,3\)，\((E,O)\) 为 \(k\) 上的椭圆曲线。则存在 \(A,B\in k\)，使 \((E,O)\) 同构于平面曲线
	\begin{equation}\label{b3:eq:H-short-weierstrass}
		Y^2Z=X^3+AXZ^2+BZ^3,\qquad O=[0:1:0],
	\end{equation}
	其中
	\[
		\Delta\coloneqq-16(4A^3+27B^2)\ne0.
	\]
	反过来，满足这个条件的平面三次方程与指定点 \([0:1:0]\) 构成一条椭圆曲线。
\end{theorem}
\begin{proof}
	由\cref{b3:cor:H-canonical-degree}，\(E\) 上的典范除子满足 \(K\sim0\)。对 \(n\ge1\)，\(K-nO\) 的次数为负，故没有非零截面；\cref{b3:thm:H-riemann-roch}给出
	\[
		\ell(nO)=n.
	\]
	特别地，\(L(O)=k\)。选取
	\[
		x\in L(2O)\setminus L(O),\qquad
		y\in L(3O)\setminus L(2O).
	\]
	它们在 \(O\) 的极点阶数分别恰为二、三，在其余点均正则。于是 \(1,x,y,x^2,xy,x^3\) 的极点阶数依次为 \(0,2,3,4,5,6\)，彼此不同，因而线性无关。它们与 \(y^2\) 都属于六维空间 \(L(6O)\)，所以有一条线性关系。关系中 \(y^2\) 与 \(x^3\) 的系数均非零，否则剩余项的最高极点阶数不能相消。适当缩放 \(y\) 后，可写成
	\[
		y^2+a_1xy+a_3y=x^3+a_2x^2+a_4x+a_6.
	\]
	代数闭性保证所需的非零标量有平方根。因二在 \(k\) 中可逆，把 \(y\) 换成 \(y+(a_1x+a_3)/2\)，可消去左端的一次 \(y\) 项；再因三可逆，平移 \(x\) 可消去右端的二次项。这些变换不改变 \(x,y\) 的极点阶数，所得关系即为 \(y^2=x^3+Ax+B\)。

	还须说明所得方程确实表示原来的光滑曲线。\Cref{b3:prop:H-high-degree-linear-systems}适用于次数为三的除子 \(3O\)，给出由完整线性系 \(|3O|\) 定义的闭嵌入
	\[
		E\longrightarrow\mathbb P^2_k,
		\qquad P\longmapsto[x(P):y(P):1]\quad(P\ne O).
	\]
	在 \(O\) 附近用局部统一化元 \(t\) 同乘 \(t^3\)，三个截面变为 \(t^3x,t^3y,t^3\)，前后两项在 \(O\) 消失，中间一项为单位，所以 \(O\) 的像为 \([0:1:0]\)。这里闭嵌入所用的输入是 \(3O\) 分离不同点及同一点的一阶切向，而不仅是三个函数线性无关。三次多项式 \(x^3+Ax+B\) 在射影直线的无穷远处有奇数阶极点，故在 \(k(x)\) 中不是平方。因此 \(y^2-x^3-Ax-B\) 不可约；闭嵌入的像遂恰为\eqref{b3:eq:H-short-weierstrass}的整个三次曲线。

	设 \(f(x)=x^3+Ax+B\)。无穷远处只有 \(O\)，且齐次方程对 \(Z\) 的偏导数在 \(O\) 为一。仿射部分出现奇点，当且仅当
	\[
		y=0,\qquad f(x)=f'(x)=0.
	\]
	三次多项式的判别式为 \(-4A^3-27B^2\)，故光滑性等价于 \(\Delta\ne0\)。也可直接核验：重根 \(r\) 必满足 \(A=-3r^2\)、\(B=2r^3\)；反向在 \(A\ne0\) 时取 \(r=-3B/(2A)\)，在 \(A=0\) 时取 \(r=0\)。原来的曲线光滑，因而它的方程必满足这一条件。

	最后证明反向的亏格断言。对满足 \(\Delta\ne0\) 的方程，有理微分
	\[
		\omega=\frac{\mathrm{d}x}{2y}
	\]
	在仿射开集 \(y\ne0\) 上是局部微分模的生成元；在 \(y=0\) 处，由 \(f'(x)\ne0\) 及 \(2y\,\mathrm{d}y=f'(x)\,\mathrm{d}x\)，它又写成 \(\mathrm{d}y/f'(x)\)，仍是生成元。在 \(O\) 附近取 \(u=X/Y\)、\(v=Z/Y\)，方程为
	\[
		v=u^3+Auv^2+Bv^3.
	\]
	其中 \(u\) 是统一化元，\(1-2Auv-3Bv^2\) 为单位。微分上式并代入 \(x=u/v\)、\(y=1/v\)，得到
	\[
		\omega=-\frac{\mathrm{d}u}{1-2Auv-3Bv^2}.
	\]
	所以 \(\omega\) 处处无零点、无极点，其除子为零。再由\cref{b3:cor:H-canonical-degree}的 \(\deg K=2g-2\)，得到 \(g=1\)。
\end{proof}

上述方程称为\term[duan-Weierstrass-fangcheng]{短 Weierstrass 方程}[short Weierstrass equation]，\(\Delta\) 称为这一方程的判别式。证明中二与三的可逆性分别用于配方和平移消去二次项；特征二、三的椭圆曲线仍有 Weierstrass 方程，但不能统一采用这个短式。沿用证明中的 \(u,v\)，还有 \(\operatorname{ord}_O(v)=3\)，从而 \(\operatorname{ord}_O(x)=-2\)、\(\operatorname{ord}_O(y)=-3\)，与最初的截面选择一致。

\begin{example}\label{b3:exm:H-weierstrass-divisors}
	设 \(k\) 为代数闭域，\(\operatorname{char}k\ne2,3\)，考虑 \(E:y^2=x^3-x\)。这里 \(\Delta=64\)，故曲线光滑。记 \(T_a=(a,0)\)，其中 \(a=-1,0,1\)。在每个 \(T_a\) 处，\(y\) 是统一化元，\(x-a\) 的阶数为二；在 \(O\) 处，\(x,y\) 的极点阶数为二、三。因此
	\[
		\operatorname{div}(x)=2T_0-2O,
		\qquad
		\operatorname{div}(y)=T_{-1}+T_0+T_1-3O.
	\]
	这些是函数除子的等式。下一小节将把它们转写为点的群运算关系。
\end{example}

\subsection{割线、切线与群律}
\label{b3:subsec:H04:02}

一条直线与光滑三次曲线相交，计重数得到三个点。用这条直线作函数除子，所得关系属于除子类群；这才是割线与群运算之间的联系。为区别点的运算和除子的形式相加，暂把点的运算记作 \(P\oplus Q\)。

\begin{theorem}[椭圆曲线的群律]\label{b3:thm:H-elliptic-group-law}
	设 \(k\) 为代数闭域，\((E,O)\) 为 \(k\) 上的椭圆曲线。映射
	\[
		\iota:E(k)\longrightarrow\operatorname{Pic}^0(E),
		\qquad P\longmapsto[P-O]
	\]
	为双射。存在唯一以 \(O\) 为单位元的交换群运算 \(\oplus\)，使 \(\iota\) 为群同构；具体地，
	\[
		P\oplus Q=R\quad\Longleftrightarrow\quad P+Q\sim R+O.
	\]
	若另有 \(\operatorname{char}k\ne2,3\)，并采用\eqref{b3:eq:H-short-weierstrass}，则任一直线的交点除子 \(P+Q+R\) 满足
	\[
		P\oplus Q\oplus R=O,
	\]
	其中允许重复点，并按交点重数计数。
\end{theorem}
\begin{proof}
	对任意次数零除子类 \([D]\)，\(D+O\) 的次数为一。由\cref{b3:prop:H-genus-one-sections}，它的类中存在唯一点 \(P\)，于是 \([D]=[P-O]\)。这同时证明 \(\iota\) 的满射性和单射性。把 \(\operatorname{Pic}^0(E)\) 的加法经此双射传回 \(E(k)\)，便得到所述运算。因为目标是交换群，结合律、交换律与逆元的存在均随之成立；\([O-O]=0\) 给出单位元。定义等式
	\[
		[P-O]+[Q-O]=[R-O]
	\]
	又正好等价于 \(P+Q\sim R+O\)。

	在上述三次嵌入中，直线 \(Z=0\) 截出的除子为 \(3O\)。任意另一条直线 \(\ell=0\) 与它的比 \(\ell/Z\) 是曲线上的有理函数，其主除子为 \(D_\ell-3O\)。所以 \(D_\ell\sim3O\)，且 \(\deg D_\ell=3\)。写作 \(D_\ell=P+Q+R\) 后，便有 \([P-O]+[Q-O]+[R-O]=0\)。这里用的是截面零除子的次数，不需要把实平面上的图形当作结合律的证明。
\end{proof}

对于仿射点 \(P=(a,b)\)，竖直直线 \(X-aZ=0\) 截出 \(P+(a,-b)+O\)，计重数的说法也包括 \(b=0\)。因此
\[
	-P=(a,-b),\qquad -O=O.
\]
两点确定的割线若在第三点 \(R\) 相交，则 \(P\oplus Q=-R\)；两点重合时用切线，并把这个点计至少两次。特别地，\(O\) 的切线 \(Z=0\) 与曲线的交点除子为 \(3O\)，所以 \(O\) 是拐点。

\begin{proposition}[割线与切线公式]\label{b3:prop:H-chord-tangent-formulas}
	设 \(k\) 为代数闭域，\(\operatorname{char}k\ne2,3\)，\(E:y^2=x^3+Ax+B\) 满足 \(\Delta\ne0\)，原点为 \(O\)。对仿射点 \(P=(x_1,y_1)\)、\(Q=(x_2,y_2)\)，若 \(x_1\ne x_2\)，令
	\[
		m=\frac{y_2-y_1}{x_2-x_1};
	\]
	若 \(P=Q\) 且 \(y_1\ne0\)，令
	\[
		m=\frac{3x_1^2+A}{2y_1}.
	\]
	在这两种情形下，\(P\oplus Q=(x_3,y_3)\)，其中
	\begin{equation}\label{b3:eq:H-addition-formulas}
		x_3=m^2-x_1-x_2,\qquad
		y_3=m(x_1-x_3)-y_1.
	\end{equation}
	其余情形由 \(P\oplus O=P\) 及 \(P\oplus(-P)=O\) 决定；特别地，\(y_1=0\) 时 \(P\oplus P=O\)。
\end{proposition}
\begin{proof}
	非竖直的割线或切线写成 \(y=m(x-x_1)+y_1\)。代入曲线方程，得到首一三次方程
	\[
		x^3-m^2x^2+\text{一次及常数项}=0.
	\]
	它的三个根是 \(x_1,x_2,x_R\)，切线时前两个根相同。因此 \(x_R=m^2-x_1-x_2\)。第三交点的纵坐标为 \(m(x_R-x_1)+y_1\)，取逆元即得 \(y_3\)。切线斜率由 \(2y_1\,\mathrm{d}y=(3x_1^2+A)\,\mathrm{d}x\) 给出。若 \(x_1=x_2\)，方程使 \(y_2=y_1\) 或 \(y_2=-y_1\)；除已列出的非竖直切线情形外，均落入竖直直线的情形，故已穷尽。
\end{proof}

\begin{example}\label{b3:exm:H-point-addition}
	在 \(k=\overline{\mathbb Q}\) 上取 \(E:y^2=x^3-x+1\)，有 \(\Delta=-16\cdot23\ne0\)。点 \(P=(0,1)\)、\(Q=(1,1)\) 之间的割线为 \(y=1\)，第三交点为 \((-1,1)\)，所以 \(P\oplus Q=(-1,-1)\)。在 \(P\) 处的切线斜率为 \(-1/2\)，从而
	\[
		[2]P\coloneqq P\oplus P
		=\left(\frac14,-\frac78\right).
	\]
	代回右端可核验 \((1/4)^3-1/4+1=49/64\)。两种计算分别使用割线和切线，所得运算的结合律已经由除子类群保证。
\end{example}

点集上有群运算，还不足以说明它与代数几何相容。特别是\eqref{b3:eq:H-addition-formulas}的分母消失时，不能把分式直接当成已定义的坐标。

\begin{proposition}\label{b3:prop:H-elliptic-group-morphisms}
	设 \(k\) 为代数闭域，\(\operatorname{char}k\ne2,3\)，\((E,O)\) 为 \(k\) 上的椭圆曲线。取逆映射 \(E\rightarrow E\)、加法映射 \(\mu:E\times E\rightarrow E\) 以及每个平移 \(\tau_A:P\mapsto A\oplus P\) 均为 \(k\) 态射；平移是曲线自同构。
\end{proposition}
\begin{proof}
	先固定 \(A\in E(k)\)。由\cref{b3:prop:H-genus-one-sections}和\cref{b3:prop:H-high-degree-linear-systems}，次数二的除子 \(A+O\) 有两个线性无关截面，且无基点，因而给出态射 \(f_A:E\rightarrow\mathbb P^1_k\)，纤维除子属于 \(|A+O|\)，次数为二。由\cref{b3:prop:H-curve-morphism}，它是二次有限态射。由于特征不为二，函数域扩张是可分二次扩张，有唯一非平凡自同构。应用\cref{b3:lem:H-rational-map-extension}于对应的有理自映射及其逆，得到 \(E\) 的自同构 \(\sigma_A\)：两个复合在函数域上为恒等，由延拓唯一性处处为恒等。

	在不分歧纤维中，\(\sigma_A\) 交换两个点；在分歧纤维中，它固定唯一的点，纤维为该点的二倍。因此对每个 \(P\) 都有
	\[
		P+\sigma_A(P)\sim A+O,
		\qquad \sigma_A(P)=A\oplus(-P).
	\]
	于是 \(\sigma_O\) 就是取逆映射，而 \(\tau_A=\sigma_A\circ\sigma_O\)。这证明了取逆和平移的代数性。

	在 \(E\times E\) 的非空开集 \(U\) 上，要求两个点均为仿射点且横坐标不同；\eqref{b3:eq:H-addition-formulas}在那里定义态射 \(\mu_U:U\rightarrow E\)。对任意闭点 \((P,Q)\)，可选 \(A,B\in E(k)\)，使 \((P\oplus A,Q\oplus B)\in U\)。例如先选 \((P_0,Q_0)\in U\)，再令 \(A=P_0\oplus(-P)\)、\(B=Q_0\oplus(-Q)\)。于是
	\[
		\tau_{-(A\oplus B)}\circ\mu_U\circ(\tau_A\times\tau_B)
	\]
	在包含 \((P,Q)\) 的开邻域上是态射，并在每个点处都等于抽象群运算 \(\oplus\)。这些邻域覆盖整个 \(E\times E\)：其补集若非空，就是有限型 \(k\) 概形中的非空闭集，必有闭点，与构造矛盾。这些局部态射在重叠处相同，遂黏合成 \(\mu\)。这里相同也可按态射核验：源为约化 \(k\) 概形，目标分离，闭点上的相等迫使态射相等。其结合律和单位元等恒等式同样由闭点上的群公理得到。
\end{proof}

这个证明用平移把非空开集上的加法公式移到任意一对点的邻域，补上了分母消失处的定义。二次线性系和这一延拓方法见\cite[Chapter IV, Lemma 4.2, Proposition 4.8, pp. 318, 321--322]{Hartshorne1977}。

\subsection{j 不变量与同源}
\label{b3:subsec:H04:03}

短 Weierstrass 方程并不唯一。改变 \(L(2O)\)、\(L(3O)\) 的基会改变系数；因此，要比较两条曲线，须先确定哪些系数变化只是在改写同一条曲线。

\begin{proposition}[短方程的同构与 \(j\) 不变量]\label{b3:prop:H-j-invariant}
	设 \(k\) 为代数闭域，\(\operatorname{char}k\ne2,3\)，两条椭圆曲线分别为
	\[
		E:y^2=x^3+Ax+B,\qquad
		E':y'^2=x'^3+A'x'+B',
	\]
	两式判别式均非零，原点均为无穷远点。每个保持原点的同构 \(E\rightarrow E'\) 恰有形式
	\[
		(x,y)\longmapsto(u^2x,u^3y),\qquad
		A'=u^4A,\quad B'=u^6B,
		\quad u\in k^\times.
	\]
	数值
	\begin{equation}\label{b3:eq:H-j-invariant}
		j(E)\coloneqq1728\,\frac{4A^3}{4A^3+27B^2}
	\end{equation}
	只依赖椭圆曲线的同构类，称为\term[j-bubianliang]{\(j\) 不变量}[j-invariant]。两条这样的椭圆曲线同构，当且仅当 \(j(E)=j(E')\)。
\end{proposition}
\begin{proof}
	同构保持 \(O\) 处的极点阶数，故它在函数上的拉回必有形式
	\[
		x'=ax+r,\qquad y'=by+sx+t,\qquad a,b\ne0.
	\]
	把它代入 \(E'\) 的方程，并用 \(y^2=x^3+Ax+B\) 约化。因为 \(1,y\) 在 \(k(x)\) 上线性无关，比较一次 \(y\) 项得 \(s=t=0\)；比较 \(x^2\) 项得 \(r=0\)；比较 \(x^3\) 项得 \(b^2=a^3\)。令 \(u=b/a\)，便有 \(u^2=a\)、\(u^3=b\)。比较其余两项即得 \(A'=u^4A\)、\(B'=u^6B\)。反向的缩放及其逆显然给出保持原点的同构。

	缩放使 \(A^3,B^2\) 同乘 \(u^{12}\)，所以\eqref{b3:eq:H-j-invariant}不变。现在设两条曲线的 \(j\) 相同。若 \(j=0\)，则 \(A=A'=0\)，而 \(B,B'\ne0\)，取 \(u^6=B'/B\) 即可。若 \(j=1728\)，则 \(B=B'=0\)，而 \(A,A'\ne0\)，取 \(u^4=A'/A\) 即可。其余情形四个系数均非零。置 \(r=A'/A\)、\(s=B'/B\)，相同的 \(j\) 给出 \(r^3=s^2\)。取 \(a=s/r\)，则 \(a^2=r\)、\(a^3=s\)；再取 \(u^2=a\)，便同时满足 \(u^4=r\)、\(u^6=s\)。所需根的存在用到了 \(k\) 的代数闭性。
\end{proof}

例如，\(y^2=x^3-x\) 的 \(j\) 为 \(1728\)，\(y^2=x^3+1\) 的 \(j\) 为零。若只比较不带原点的曲线同构，一样得到这个分类：给定同构 \(f:E\rightarrow E'\)，再复合目标上的平移 \(\tau_{-f(O)}\)，就得到保持原点的同构。这里 \(j\) 的分类范围是代数闭域上的光滑曲线同构；它不是同源的分类。Hartshorne 采用 Legendre 方程给出同一不变量的分类\cite[Chapter IV, Theorem 4.1, pp. 317--320]{Hartshorne1977}。

\begin{definition}\label{b3:def:H-isogeny}
	设 \(k\) 为代数闭域，\(\operatorname{char}k\ne2,3\)，\((E,O)\)、\((E',O')\) 为 \(k\) 上的椭圆曲线。一个非恒定 \(k\) 态射 \(\varphi:E\rightarrow E'\)，若满足
	\[
		\varphi(O)=O',\qquad
		\varphi(P\oplus Q)=\varphi(P)\oplus\varphi(Q)
		\quad(P,Q\in E(k)),
	\]
	则称为一个\term[tongyuan]{同源}[isogeny]。其次数定义为
	\[
		\deg\varphi=[k(E):\varphi^*k(E')],
	\]
	包括函数域扩张的不可分次数。若该扩张可分，则称 \(\varphi\) 为可分同源。记 \(\ker\varphi(k)=\{\,P\in E(k)\mid\varphi(P)=O'\,\}\)，这是有限的点群；概形纤维 \(\varphi^{-1}(O')\) 还保留各点的重数。
\end{definition}

在这个定义中，群条件其实由保持原点的非恒定曲线态射自动满足。下面给出所需的除子理由，而不把它作为额外约定。

\begin{lemma}\label{b3:lem:H-origin-morphism-homomorphism}
	设 \(k\) 为代数闭域，\(\operatorname{char}k\ne2,3\)，\((E,O)\)、\((E',O')\) 为椭圆曲线。每个满足 \(f(O)=O'\) 的非恒定 \(k\) 态射 \(f:E\rightarrow E'\) 都是同源。
\end{lemma}
\begin{proof}
	由\cref{b3:prop:H-curve-morphism}，\(f\) 有限且平坦。将点除子 \(P\) 推为 \(f(P)\)，并线性延拓，得到 \(f_*:\operatorname{Div}(E)\rightarrow\operatorname{Div}(E')\)。这个推送保持线性等价，因为对 \(h\in k(E)^\times\)，有
	\[
		f_*\operatorname{div}(h)
		=\operatorname{div}\bigl(\operatorname{Norm}_{k(E)/f^*k(E')}(h)\bigr).
	\]
	说明这个等式的局部依据：在目标点的 DVR 上，有限平坦代数是有限秩自由模，范数是乘 \(h\) 的线性映射的行列式。对局部正则的非零 \(h\)，在 DVR 上逐次选取赋值最小的非零矩阵项作主元，以可逆行列变换消去同行同列的其余项，便把矩阵化为对角形。对角元阶数之和同时等于行列式的赋值与余核长度；按源中的各点分解，这一长度是各 \(\operatorname{ord}_P(h)\) 之和，因剩余域均为 \(k\)。对一般有理函数清去分母并相减，即得所写等式。这也适用于不可分扩张。

	若 \(P\oplus Q=R\)，则 \(P+Q\sim R+O\)。推送后得到
	\[
		f(P)+f(Q)\sim f(R)+O',
	\]
	由\cref{b3:thm:H-elliptic-group-law}，这恰好是 \(f(P)\oplus f(Q)=f(R)\)。
\end{proof}

这一结论也见\cite[Chapter IV, Lemma 4.9, p. 322]{Hartshorne1977}。同源的纤维是核的平移，但一个纤维的点数和它的次数仍可能不同。

\begin{proposition}[同源的次数与核]\label{b3:prop:H-isogeny-degree}
	设 \(k\) 为代数闭域，\(\operatorname{char}k\ne2,3\)，\(\varphi:(E,O)\rightarrow(E',O')\) 为同源，次数为 \(d\)。则 \(\varphi\) 有限、平坦且满射。设 \(e\) 为它在 \(O\) 处的分歧指数，则在每个闭点的分歧指数都为 \(e\)，并且
	\[
		d=e\,\#\ker\varphi(k).
	\]
	下列条件等价：\(\varphi\) 可分；\(e=1\)；\(\varphi\) 在每个点都不分歧；\(\#\ker\varphi(k)=d\)。
\end{proposition}
\begin{proof}
	前三项性质由\cref{b3:prop:H-curve-morphism}给出。对 \(P\in E(k)\)，群同态性质给出交换图
	\[
		\begin{tikzcd}
			E \arrow[r,"\tau_P"] \arrow[d,"\varphi"'] & E \arrow[d,"\varphi"] \\
			E' \arrow[r,"\tau_{\varphi(P)}"'] & E' .
		\end{tikzcd}
	\]
	水平箭头均为自同构，所以 \(\varphi\) 在 \(P\) 与 \(O\) 处的局部参数阶数相同。这说明各闭点分歧指数相同。对原点纤维应用\cref{b3:prop:H-curve-morphism}的纤维次数公式，得 \(d=\sum_{P\in\ker\varphi(k)}e=e\,\#\ker\varphi(k)\)。

	令 \(L=\varphi^*k(E')\)、\(K=k(E)\)。两条曲线光滑，故 \(K\otimes_L\Omega_{L/k}\) 与 \(\Omega_{K/k}\) 都是一维 \(K\) 空间。\Cref{b3:thm:differential-transitivity}给出正合列
	\[
		K\otimes_L\Omega_{L/k}\longrightarrow\Omega_{K/k}
		\longrightarrow\Omega_{K/L}\longrightarrow0.
	\]
	第一条映射非零，恰好等价于 \(\Omega_{K/L}=0\)；由\cref{b3:thm:finite-field-differentials-separable}，这又等价于 \(K/L\) 可分。

	若 \(\varphi\) 可分，函数域处的微分非零，故在某个非空开集上的切映射非零。对局部统一化元，切映射非零恰好表示拉回有非零一次项，也就是分歧指数为一。平移使不分歧性传到全部点，所以 \(e=1\)。反过来，若 \(e=1\)，\(O\) 处的切映射非零，函数域处的微分便非零，由刚证的判据得到可分性。其余等价关系随已证公式得到。曲线态射的这一微分计算也见\cite[Chapter IV, Propositions 2.1--2.2, pp. 299--300]{Hartshorne1977}。
\end{proof}

\begin{example}[倍点同源]\label{b3:exm:H-doubling-isogeny}
	设 \(k\) 为代数闭域，\(\operatorname{char}k\ne2,3\)，\(E:y^2=x^3+Ax+B\) 满足 \(\Delta\ne0\)。反复复合加法和取逆，得到整数倍映射 \([n]:E\rightarrow E\)。它在点上是交换群同态。在 \((O,O)\) 处，加法的微分限制到两个坐标轴都是恒等，所以
	\[
		\mathrm{d}\mu_{(O,O)}(v,w)=v+w,
		\qquad \mathrm{d}[n]_O=n\,\operatorname{id}_{T_OE}.
	\]
	特别地，\([2]\) 的微分非零，故它非恒定且可分。又因 \([2]P=O\) 等价于 \(P=-P\)，其点核恰为
	\[
		E[2](k)=\{\,O,(r_1,0),(r_2,0),(r_3,0)\,\},
	\]
	其中 \(r_1,r_2,r_3\) 为 \(x^3+Ax+B\) 的三个不同根。于是\cref{b3:prop:H-isogeny-degree}给出 \(\deg[2]=4\)，而且 \(E[2](k)\cong(\mathbb Z/2\mathbb Z)^2\)。在 \(y\ne0\) 的仿射部分，\([2]\) 的坐标由\eqref{b3:eq:H-addition-formulas}取 \(P=Q\) 算出；在其余点的值由已经建立的态射与群律决定。
\end{example}

\symidx{\([n]\)：椭圆曲线的整数倍映射}

若正特征整除 \(n\)，上面的切映射为零，不能据此认定 \([n]\) 恒定，也不能直接用核的点数计算次数。在特征五的三次曲线上，可以直接算出一个不可分同源。

\begin{example}\label{b3:exm:H-inseparable-isogeny}
	取 \(k=\overline{\mathbb F}_5\)，\(E:Y^2Z=X^3-XZ^2\)，原点为 \(O=[0:1:0]\)。因方程系数属于 \(\mathbb F_5\)，坐标式
	\[
		F:E\longrightarrow E,
		\qquad [X:Y:Z]\longmapsto[X^5:Y^5:Z^5]
	\]
	定义一个保持原点的非恒定 \(k\) 态射，故由\cref{b3:lem:H-origin-morphism-homomorphism}，它是同源。它在 \(k\) 点上为双射，特别地 \(\ker F(k)=\{O\}\)。然而在 \(O\) 的统一化元 \(u=X/Y\) 上，\(F^*u=u^5\)，所以 \(e=5\)，进而 \(\deg F=5\)。所有坐标微分均为零，与不可分性相符。这个 \(k\) 态射在坐标环上固定 \(k\) 的元素；它不同于把常数也取五次方的绝对 Frobenius 环映射。
\end{example}

三次方程可以用来核验光滑性、计算群运算和比较同构类；同源则比较保持群结构的有限态射，其次数必须连同可分性一起判断。一般整数倍映射的次数、挠点的完整结构及椭圆曲线的算术理论需要进一步的工具，这里只使用已经计算的倍点与五次方态射。

% BEGIN APPENDIX OUTLINE SYNC
% * 状态：初稿已完成；新增范围尚未正式审核。
% #### H.4.1 Weierstrass 方程与光滑平面三次曲线
% #### H.4.2 割线、切线与群律
% #### H.4.3 j 不变量与同源
% ---
% END APPENDIX OUTLINE SYNC


\begin{chaptersummary}
	在光滑曲线的每个闭点，局部环的 DVR 结构使零极点阶数成为整数。射影性让非恒定函数给出到射影直线的有限满射，因此零点和极点的总重数相等。关联层 \(\mathcal O_C(D)\) 把除子条件转为线丛；线性等价对应线丛同构，次数则成为线丛的数值。

	\ProseParagraphBreak
	对线丛逐次增加一个点，Euler 特征每次增加一。Serre 对偶把一次上同调识别为典范线丛所控制的截面空间的对偶，由此得到截面维数公式及 \(\deg K_C=2g-2\)。射影直线的 Laurent 指数计算同时显示正次数的截面与负次数的一次上同调；亏格一曲线的零次线丛则说明，次数相同仍可能有不同的上同调。

	\ProseParagraphBreak
	选定亏格一曲线上的原点后，每个次数零的除子类有唯一的形式 \([P-O]\)。群律的结合性由 Picard 群传递到点集，割线与切线把它变成三次方程上的计算。短 Weierstrass 模型的 \(j\) 不变量描述代数闭域上的同构类；同源则比较不同的曲线与群结构。正特征下，纤维长度和点的个数可能不同，必须保留不可分次数。
\end{chaptersummary}

\BookExercises

\begin{exercise}\label{b3:ex:H-rational-double-cover}
	设 \(k\) 为特征不为 \(2\) 的代数闭域。考虑 \(\varphi:\mathbb P_k^1\rightarrow\mathbb P_k^1\)，仿射表达式为 \(u=t+t^{-1}\)。写出其齐次表达式，计算次数以及 \(u=\infty,2,-2\) 三个纤维的除子，并解释为何纤维点数会变化。
\end{exercise}
\begin{solution}
	在坐标 \([T:S]\) 下，映射为 \([T^2+S^2:TS]\)。两个齐次式没有公共零点：若 \(TS=0\)，另一个式子仍非零。因此它是全局态射。函数 \(u=(t^2+1)/t\) 在 \(0,\infty\) 各有一阶极点，故次数为二。又有
	\[
		u-2=\frac{(t-1)^2}{t},\qquad u+2=\frac{(t+1)^2}{t}.
	\]
	于是三个纤维分别为 \([0]+[\infty]\)、\(2[1]\)、\(2[-1]\)。特征不为二使 \(1,-1\) 和 \(2,-2\) 分别不同。一般纤维由二次式 \(t^2-at+1\) 给出，当 \(a^2\ne4\) 时有两个不同的根；在 \(a=\pm2\) 时两根重合，但纤维总重数仍为二。
\end{solution}

\begin{exercise}\label{b3:ex:H-prescribed-poles}
	设 \(k\) 为代数闭域，在 \(\mathbb P_k^1\) 上取 \(D=2[0]+[1]-[\infty]\)。求 \(L(D)\) 的一组基，并在无穷远的局部坐标中核对每个基向量满足截面条件。
\end{exercise}
\begin{solution}
	令 \(h=t^2(t-1)\)，则 \(\operatorname{div}(h)=D-2[\infty]\)。乘以 \(h\) 将 \(L(D)\) 同构到 \(L(2[\infty])\)。后者由次数至多二的多项式组成，所以可取基
	\[
		\frac1{t^2(t-1)},\qquad \frac1{t(t-1)},\qquad \frac1{t-1}.
	\]
	令 \(s=t^{-1}\)，这三个函数分别是 \(s^3/(1-s)\)、\(s^2/(1-s)\)、\(s/(1-s)\)。它们在 \(\infty\) 的阶数分别为 \(3,2,1\)，都至少为一，恰满足 \(D\) 在该点系数为 \(-1\) 的要求。在 \(0\) 的极点阶数至多二，在 \(1\) 至多一，其他点无极点，故均属于 \(L(D)\)。乘以 \(h\) 后得到 \(1,t,t^2\)，同时证明了生成性与线性无关性。
\end{solution}

\begin{exercise}\label{b3:ex:H-cusp-rational-map}
	设 \(k\) 为代数闭域，\(X=\operatorname{Spec}k[t^2,t^3]\)，\(P\) 为极大理想 \((t^2,t^3)\) 所给尖点。证明从 \(X\) 到 \(\mathbb P_k^1\) 的有理映射 \(t\mapsto[t:1]\) 不能跨过 \(P\) 延拓。说明它在正规化 \(\operatorname{Spec}k[t]\) 上为何可以延拓。
\end{exercise}
\begin{solution}
	记 \(R=k[t^2,t^3]_{(t^2,t^3)}\)。若延拓在 \(P\) 的像是有限点，则坐标 \(t\) 必须属于 \(R\)。写 \(t=f/g\)，其中 \(f,g\in k[t^2,t^3]\)、\(g(0)\ne0\)，则 \(tg=f\) 的一次项系数为 \(g(0)\ne0\)，而 \(f\) 没有一次项，矛盾。若像为无穷远点，则 \(t^{-1}\in R\)，但 \(R\subset k[t]_{(t)}\)，右边的每个元素在 \(0\) 处阶数非负，\(t^{-1}\) 的阶数为 \(-1\)，仍矛盾。正规化的坐标环已经含 \(t\)，所以 \([t:1]\) 在整个仿射直线上定义。这说明射影目标本身不能弥补源的非正规局部环。
\end{solution}

\begin{exercise}\label{b3:ex:H-residue-degree-weight}
	本题把底域改为 \(\mathbb R\)。在 \(\mathbb P_{\mathbb R}^1\) 上令 \(P\) 为不可约多项式 \(t^2+1\) 所给闭点。计算 \(\operatorname{div}(t^2+1)\)，并分别用系数直接相加、按剩余域次数加权计算其次数。
\end{exercise}
\begin{solution}
	\(\mathbb R[t]_{(t^2+1)}\) 是 DVR，\(t^2+1\) 为统一化元，故在 \(P\) 的阶数是一。无穷远处写成 \(s^{-2}(1+s^2)\)，阶数为 \(-2\)。因此主除子为 \(P-2[\infty]\)。不加权的系数和为 \(-1\)，没有满足次数零公式。但 \(\kappa(P)=\mathbb R[t]/(t^2+1)\cong\mathbb C\)，剩余域次数为二，而 \(\infty\) 为实有理点。因此正确的加权次数为 \(1\cdot2-2\cdot1=0\)。在代数闭域上的公式中，各权重都等于一，才可以省略。
\end{solution}

\begin{exercise}\label{b3:ex:H-fixed-zeros-linear-system}
	设 \(k\) 为代数闭域，在 \(\mathbb P_k^1\) 上用 \(t\) 表示有限图的坐标，令
	\[
		D=4P_\infty-P_0-2P_1.
	\]
	求 \(L(D)\) 的一组基。所有函数都在 \(P_0,P_1\) 为零，是否意味着 \(|D|\) 有基点？用这组基写出线性系给出的态射，并在 \(P_0,P_1,P_\infty\) 处用线丛局部基核验。
\end{exercise}
\begin{solution}
	截面是次数至多四、被 \(t(t-1)^2\) 整除的多项式。因此令 \(h=t(t-1)^2\)，就有
	\[
		L(D)=\operatorname{span}_k\{h,th\}.
	\]
	在 \(P_0\) 处，\(D\) 的系数为 \(-1\)，线丛局部基可取 \(t\)。两个截面的系数为 \((t-1)^2,t(t-1)^2\)，第一个在 \(P_0\) 不为零。在 \(P_1\) 处，局部基可取 \((t-1)^2\)，系数为 \(t,t^2\)，也没有共同零点。

	在无穷远处取 \(u=t^{-1}\)，局部基为 \(t^4\)。两个系数分别为
	\[
		\frac{h}{t^4}=u(1-u)^2,\qquad
		\frac{th}{t^4}=(1-u)^2.
	\]
	第二个在 \(u=0\) 不为零。其他有限点处，局部基为 \(1\)，而 \(h\ne0\)。所以线性系没有基点。在一般点，态射是 \([h:th]=[1:t]\)；无穷远图上的表达式为 \([u:1]\)。它是射影直线的一个同构。函数本身的零点已被除子中的负系数抵消，必须在局部基下判断截面是否消失。
\end{solution}

\begin{exercise}\label{b3:ex:H-genus-one-degree-two}
	设 \(k\) 为代数闭域，\(C\) 为亏格一的光滑射影整曲线，\(D\) 为次数二的除子。证明 \(|D|\) 定义到 \(\mathbb P_k^1\) 的次数二态射，因此不能给出闭嵌入。证明对每个闭点 \(P\)，都有唯一闭点 \(Q\) 满足 \(P+Q\sim D\)，且点集映射 \(\iota_D(P)=Q\) 满足 \(\iota_D^2=\mathrm{id}\)。
\end{exercise}
\begin{solution}
	由\cref{b3:prop:H-genus-one-sections}，\(\ell(D)=2\)；由\cref{b3:prop:H-high-degree-linear-systems}，\(|D|\) 无基点。因此它定义态射 \(\varphi:C\rightarrow\mathbb P^1\)，且 \(\varphi^*\mathcal O(1)\cong\mathcal O_C(D)\)。两独立截面的比不是常数，所以 \(\varphi\) 非恒定。一个目标点的拉回是线性系中的次数二有效除子，\cref{b3:prop:H-curve-morphism}于是给出 \(\deg\varphi=2\)。若为闭嵌入，其像作为一维闭整子概形只能是 \(\mathbb P^1\)，次数应为一，矛盾。

	\(D-P\) 次数为一，故其除子类含有唯一闭点 \(Q\)。这正是 \(P+Q\sim D\)。交换 \(P,Q\) 后，\(P\) 又是 \(D-Q\) 所属类中的闭点；唯一性给出 \(\iota_D(Q)=P\)。该结论允许 \(P=Q\)，不需要假定特征不为二，也不以纤维中不同点的个数替代次数。
\end{solution}

\begin{exercise}\label{b3:ex:H-base-point-bound-sharp}
	设 \(k\) 为代数闭域，\(C\) 为亏格 \(g\ge1\) 的光滑射影整曲线，\(K_C\) 为典范除子，\(P\) 为闭点。证明 \(\ell(K_C+P)=\ell(K_C)=g\)，且 \(P\) 是 \(|K_C+P|\) 的基点。说明这为何表明统一的无基点次数界 \(2g\) 不能改为 \(2g-1\)。
\end{exercise}
\begin{solution}
	Riemann–Roch 用于 \(K_C+P\)，其修正项是 \(\ell(-P)=0\)，故
	\[
		\ell(K_C+P)=(2g-1)+1-g=g.
	\]
	另一方面，\cref{b3:cor:H-canonical-degree}给出 \(\ell(K_C)=g\)。自然包含 \(L(K_C)\subseteq L(K_C+P)\) 的两端维数相同，故两空间相等。

	若 \(K_C\) 在 \(P\) 处的系数为 \(m\)，统一化元为 \(t\)，则两线丛的局部基为 \(t^{-m}\) 与 \(t^{-m-1}\)。来自 \(L(K_C)\) 的截面在后一组局部基下，其系数都多出一个因子 \(t\)，所以全部在 \(P\) 消失。因全部截面都来自这个子空间，\(P\) 是基点。此除子的次数恰为 \(2g-1\)，已给出不能降低统一界的例子。
\end{solution}

\begin{exercise}\label{b3:ex:H-genus-two-degree-three}
	设 \(k\) 为代数闭域，\(C\) 为亏格二的光滑射影整曲线，\(D\) 为次数三的除子。证明 \(\ell(D)=2\)，并证明闭点 \(P\) 是 \(|D|\) 的基点，当且仅当 \(D-P\sim K_C\)。进一步证明基点若存在就唯一，且去掉该固定点后的线性系是典范线性系。
\end{exercise}
\begin{solution}
	因为 \(\deg D=3>2g-2=2\)，高次数公式给出 \(\ell(D)=3+1-2=2\)。对任意次数二的除子 \(E\)，Riemann–Roch 给出
	\[
		\ell(E)=1+\ell(K_C-E).
	\]
	\(K_C-E\) 次数为零，其截面维数在 \(E\sim K_C\) 时为一，否则为零。因此 \(\ell(E)=2\) 当且仅当 \(E\sim K_C\)。取 \(E=D-P\)，再用基点条件 \(L(D-P)=L(D)\)，便得所求等价。

	若不同点 \(P,Q\) 都为基点，就有 \(D-P\sim K_C\sim D-Q\)，从而 \(P\sim Q\)。相应有理函数有一个简单极点，给出到 \(\mathbb P^1\) 的次数一态射。\cref{b3:prop:H-curve-morphism}使 \(C\cong\mathbb P^1\)，与亏格二矛盾，所以至多有一个基点。若该点为 \(P\)，则 \(L(D)=L(D-P)\)，且 \(D-P\sim K_C\)。由\cref{b3:exm:H-genus-two-canonical}，后者无基点，所以固定部分恰为 \(P\)；去掉它得到典范线性系，所给态射次数为二。
\end{solution}

\begin{exercise}\label{b3:ex:H-duality-matrix}
	设 \(k\) 为代数闭域。在 \(\mathbb P_k^1\) 上，用截面 \(t^2\in H^0(\mathcal O(2))\) 定义线丛同态
	\[
		u:\mathcal O(-5)\longrightarrow\mathcal O(-3).
	\]
	采用\cref{b3:prop:H-projective-line-duality}的局部基，省略基符号后，分别以 \([t^{-4}],[t^{-3}],[t^{-2}],[t^{-1}]\) 和 \([t^{-2}],[t^{-1}]\) 为两端一次上同调的有序基，求 \(H^1(u)\) 的矩阵。对偶的截面映射是
	\[
		u^*:H^0(\omega_{\mathbb P^1}\otimes\mathcal O(3))
		\longrightarrow H^0(\omega_{\mathbb P^1}\otimes\mathcal O(5)).
	\]
	以 \(\mathrm{d}t,t\,\mathrm{d}t\) 和 \(\mathrm{d}t,t\,\mathrm{d}t,t^2\mathrm{d}t,t^3\mathrm{d}t\) 为这两空间的有序基（仍省略逆线丛的局部基符号），求其矩阵，并用系数配对核验自然性。为何这两个矩阵表面上并非互为转置？
\end{exercise}
\begin{solution}
	乘以 \(t^2\) 将前两个上同调基向量送到目标的两个基向量，将后两个送到 \([1],[t]\)，而后两类在余核中为零。所以按列记录基向量的像时，矩阵为
	\[
		A=\begin{pmatrix}1&0&0&0\\0&1&0&0\end{pmatrix}.
	\]
	截面映射也是乘以 \(t^2\)，将 \(\mathrm{d}t,t\,\mathrm{d}t\) 送到 \(t^2\mathrm{d}t,t^3\mathrm{d}t\)。所列单项式基下的矩阵为
	\[
		B=\begin{pmatrix}0&0\\0&0\\1&0\\0&1\end{pmatrix}.
	\]
	对任意代表元 \(h\) 和截面系数 \(q\)，有
	\[
		\operatorname{coeff}_{t^{-1}}((t^2h)q)
		=\operatorname{coeff}_{t^{-1}}(h(t^2q)),
	\]
	故 \(\langle H^1(u)[h],q\rangle=\langle[h],u^*q\rangle\)，逐项给出自然性。这里 \([t^{-4}],[t^{-3}],[t^{-2}],[t^{-1}]\) 的对偶基依次是 \(t^3\mathrm{d}t,t^2\mathrm{d}t,t\,\mathrm{d}t,\mathrm{d}t\)，而目标两个上同调基向量的对偶基是 \(t\,\mathrm{d}t,\mathrm{d}t\)。题中截面基的顺序与这些对偶基相反，故写出的 \(B\) 并非 \(A^{\mathsf T}\)；改用这两组真正的对偶基，截面映射的矩阵便为 \(A^{\mathsf T}\)。
\end{solution}

\begin{exercise}\label{b3:ex:H-jet-evaluation-obstruction}
	设 \(k\) 为代数闭域，在 \(\mathbb P_k^1\) 上取 \(\mathcal L=\mathcal O(P_\infty)\)、\(Z=2P_0+P_\infty\)。在 \(P_0\) 处以局部基 \(1\) 记录截面系数的常数项和一次项，在 \(P_\infty\) 处以局部基 \(t\) 记录纤维值。将 \(H^0(\mathcal L|_Z)\) 因而识别为 \(k^3\)，写出求值映射并判断三项数据 \((A,B,C)\) 何时可实现。采用“后一张减前一张”的 Čech 微分，计算连接态射
	\[
		\delta:k^3\longrightarrow H^1(\mathcal L(-Z)),
	\]
	并计算它与全局截面 \(\mathrm{d}t/t^2\in H^0(\omega_{\mathbb P^1}(2P_0))\) 的 Serre 对偶配对。
\end{exercise}
\begin{solution}
	全局截面为 \(a+bt\)。在 \(P_0\) 处的两项系数是 \(a,b\)；无穷远处令 \(u=t^{-1}\)，则相对于基 \(t\) 的系数为 \(au+b\)，纤维值为 \(b\)。因此
	\[
		\mathrm{ev}_Z(a+bt)=(a,b,b),
	\]
	三项数据可实现当且仅当 \(B=C\)。

	这里 \(\mathcal L(-Z)=\mathcal O(-2P_0)\)。在有限图上的截面为 \(t^2k[t]\)，在无穷远图上的截面为 \(k[t^{-1}]\)，交叠处为 \(k[t,t^{-1}]\)。所以
	\[
		H^1(\mathcal O(-2P_0))
		\cong\frac{k[t,t^{-1}]}{t^2k[t]+k[t^{-1}]}
		=\operatorname{span}_k\{[t]\},
	\]
	其中 \([t]\) 表示单项式 \(t\) 的上同调类，故该空间是一维。三项数据在两图上的局部提升可取 \(A+Bt\)、\(Ct\)；它们的 Čech 差为
	\[
		Ct-(A+Bt)=(C-B)t-A.
	\]
	常数项在余核中为零，于是 \(\delta(A,B,C)=(C-B)[t]\)，其核恰与求值映射的像一致。

	微分 \(\mathrm{d}t/t^2\) 在 \(P_0\) 至多有二阶极点，而在无穷远处等于 \(-\mathrm{d}u\)，因此是题中线丛的全局截面。乘法把 \([t]\) 送到 \([t^{-1}\mathrm{d}t]\)。按\cref{b3:prop:H-projective-line-duality}的迹归一化，所求配对为 \(C-B\)。它把唯一的局部数据相容条件写成一个微分配对。
\end{solution}

\begin{exercise}\label{b3:ex:H-genus-one-critical-evaluation}
	设 \(k\) 为代数闭域，\(C\) 为亏格一的光滑射影整曲线，\(\mathcal L\) 为次数 \(r>0\) 的线丛，\(Z\) 为次数 \(r\) 的有效除子。证明求值映射
	\[
		\mathrm{ev}_Z:H^0(C,\mathcal L)\longrightarrow H^0(C,\mathcal L|_Z)
	\]
	的两端都是 \(r\) 维。证明若 \(\mathcal L\not\cong\mathcal O_C(Z)\)，它是同构；若 \(\mathcal L\cong\mathcal O_C(Z)\)，其核和余核都为一维。
\end{exercise}
\begin{solution}
	\cref{b3:prop:H-genus-one-cohomology}给出 \(h^0(\mathcal L)=r\)、\(H^1(\mathcal L)=0\)。在 \(Z\) 的每个支撑点，限制线丛是相应局部商环上的秩一自由模；这些商环的长度相加为 \(r\)，故 \(h^0(\mathcal L|_Z)=r\)。令 \(\mathcal M=\mathcal L(-Z)\)，限制短正合列给出
	\[
		\begin{aligned}
			0&\longrightarrow H^0(\mathcal M)\longrightarrow H^0(\mathcal L)
			\xrightarrow{\mathrm{ev}_Z}H^0(\mathcal L|_Z)\\
			&\longrightarrow H^1(\mathcal M)\longrightarrow0.
		\end{aligned}
	\]
	\(\mathcal M\) 次数为零。若 \(\mathcal L\not\cong\mathcal O_C(Z)\)，则 \(\mathcal M\) 非平凡，同一上同调表使它的零次、一次上同调同时为零，求值映射因而是同构。若 \(\mathcal L\cong\mathcal O_C(Z)\)，则 \(\mathcal M\cong\mathcal O_C\)，两次上同调都为一维；正合列便给出一维核和一维余核。这说明在次数恰好相等的边界上，能否同时满足全部局部条件取决于除子类。
\end{solution}

\begin{exercise}\label{b3:ex:H-pole-basis}
	设 \(k\) 为代数闭域，\(\operatorname{char}k\ne2,3\)，\(E:y^2=x^3+Ax+B\) 的判别式非零，\(O\) 为无穷远点。对每个整数 \(n\ge1\)，求 \(L(nO)\) 的一组由 \(1,x^i,x^jy\) 构成的基。进一步证明，在 \(E\setminus\{O\}\) 上正则的每个有理函数都唯一写成
	\[
		a(x)+b(x)y,\qquad a(x),b(x)\in k[x].
	\]
\end{exercise}
\begin{solution}
	\(x^i\) 在 \(O\) 的极点阶数为 \(2i\)，\(x^jy\) 的极点阶数为 \(2j+3\)。故所求基为
	\[
		\{1\}\ \cup\ \{\,x^i\mid i\ge1,\ 2i\le n\,\}
		\ \cup\ \{\,x^jy\mid j\ge0,\ 2j+3\le n\,\}.
	\]
	这些函数在其余点正则，而且非零项的极点阶数彼此不同，因此线性无关。\(n=1,2\) 时，基分别为 \(\{1\}\)、\(\{1,x\}\)。\(n\ge3\) 时，列出的函数数目为
	\[
		1+\left\lfloor\frac n2\right\rfloor
		+\left(\left\lfloor\frac{n-3}{2}\right\rfloor+1\right)=n.
	\]
	Riemann--Roch 给出 \(\ell(nO)=n\)，所以这组线性无关函数确实张成 \(L(nO)\)。

	设 \(h\in k(E)\) 在 \(E\setminus\{O\}\) 上正则。其极点除子仅可能支撑在 \(O\)，因而 \(h\in L(nO)\) 对某个足够大的 \(n\) 成立。按上述基展开，就得到 \(a(x)+b(x)y\)。若有两个这样的表达式，相减后得到有限个极点阶数互异的单项式的线性关系，只能所有系数均为零。故表达式唯一。也就是说，去掉原点后的正则函数环是 \(k[x,y]/(y^2-x^3-Ax-B)\)。
\end{solution}

\begin{exercise}\label{b3:ex:H-order-six-points}
	在 \(k=\overline{\mathbb Q}\) 上，取 \(E:y^2=x^3+1\)，原点为无穷远点，且令 \(P=(0,1)\)、\(Q=(-1,0)\)。计算 \([2]P\)、\(P\oplus Q\)，并求 \(P,Q,P\oplus Q\) 的阶。
\end{exercise}
\begin{solution}
	判别式为 \(-432\ne0\)，所以所给方程是椭圆曲线。\(P\) 处的切线斜率为零，故
	\[
		[2]P=(0,-1)=-P.
	\]
	因此 \([3]P=O\)，又 \(P\ne O\)，所以 \(P\) 的阶为三。\(Q\) 满足 \(Q=-Q\ne O\)，故其阶为二。

	连接 \(P,Q\) 的割线斜率为 \((0-1)/(-1-0)=1\)。代入割线公式，得到
	\[
		R\coloneqq P\oplus Q=(2,-3).
	\]
	可核验 \((-3)^2=2^3+1\)。群律给出 \([2]R=[2]P=-P\ne O\)，\([3]R=[3]Q=Q\ne O\)，而 \([6]R=O\)。因此 \(R\) 的阶恰为六；上述计算排除了六的每个真因子。
\end{solution}

\begin{exercise}\label{b3:ex:H-legendre-j}
	设 \(k\) 为代数闭域，\(\operatorname{char}k\ne2,3\)，\(\lambda\in k\setminus\{0,1\}\)。考虑 Legendre 方程
	\[
		E_\lambda:y^2=x(x-1)(x-\lambda),
	\]
	原点为无穷远点。
	\begin{enumerate}
		\item 平移 \(x\)，求其短 Weierstrass 方程，并算出判别式及 \(j(E_\lambda)\)。
		\item 直接核验 \(j(E_\lambda)=j(E_{1-\lambda})=j(E_{1/\lambda})\)，并求使 \(j=0\) 的 \(\lambda\)，以及 \(j(E_{-1})\)。
	\end{enumerate}
\end{exercise}
\begin{solution}
	1. 置 \(x=X+(1+\lambda)/3\)，得到 \(y^2=X^3+AX+B\)，其中
	\[
		A=-\frac{\lambda^2-\lambda+1}{3},\qquad
		B=-\frac{(\lambda+1)(2\lambda-1)(\lambda-2)}{27}.
	\]
	展开可得 \(4A^3+27B^2=-\lambda^2(\lambda-1)^2\)，所以
	\[
		\Delta=16\lambda^2(\lambda-1)^2\ne0,
		\qquad
		j(E_\lambda)=256\,\frac{(\lambda^2-\lambda+1)^3}{\lambda^2(\lambda-1)^2}.
	\]
	特别地，原方程光滑，且平移后仍以同一无穷远点为原点。

	2. 把 \(\lambda\) 换为 \(1-\lambda\)，分子中的二次多项式保持不变，分母也保持不变。把它换为 \(1/\lambda\)，分子乘 \(\lambda^{-6}\)，分母变为 \((\lambda-1)^2/\lambda^4\)，故商仍为原式。由于 \(256\) 在 \(k\) 中非零，\(j=0\) 当且仅当 \(\lambda^2-\lambda+1=0\)。这两个根互异，因为该二次式的判别式为 \(-3\ne0\)，且都不等于零或一。最后
	\[
		j(E_{-1})=256\,\frac{3^3}{(-1)^2(-2)^2}=1728.
	\]
	由正文的分类，这些相同 \(j\) 的方程在 \(k\) 上表示同构的椭圆曲线。
\end{solution}

\begin{exercise}\label{b3:ex:H-elliptic-automorphisms}
	设 \(k\) 为代数闭域，\(\operatorname{char}k\ne2,3\)，\(E:y^2=x^3+Ax+B\) 的判别式非零，原点为无穷远点。求保持原点的曲线自同构群 \(\operatorname{Aut}(E,O)\)，分别讨论 \(AB\ne0\)、\(B=0\) 及 \(A=0\) 的情形。对 \(y^2=x^3-x\) 写出一个四阶自同构，对 \(y^2=x^3+1\) 写出一个三阶自同构。
\end{exercise}
\begin{solution}
	正文的短方程同构判据表明，每个自同构唯一由 \(u\in k^\times\) 给出，满足
	\[
		(x,y)\longmapsto(u^2x,u^3y),\qquad
		u^4A=A,\quad u^6B=B.
	\]
	复合对应于 \(u\) 的乘法，且这个参数化是单射：若 \(u^2=u^3=1\)，则 \(u=1\)。若 \(AB\ne0\)，两条件等价于 \(u^2=1\)，所以群为 \(\mu_2(k)\cong\mathbb Z/2\mathbb Z\)。若 \(B=0\)，则 \(A\ne0\)，唯一条件为 \(u^4=1\)，群为 \(\mu_4(k)\cong\mathbb Z/4\mathbb Z\)。若 \(A=0\)，则 \(B\ne0\)，唯一条件为 \(u^6=1\)，群为 \(\mu_6(k)\cong\mathbb Z/6\mathbb Z\)。因为特征不为二、三，这些单位根均互异；域中有限乘法子群为循环群。

	在 \(y^2=x^3-x\) 上取 \(i^2=-1\)，用 \(u=i\) 得
	\[
		\alpha(x,y)=(-x,-iy),\qquad
		\alpha^2(x,y)=(x,-y),\qquad \alpha^4=\operatorname{id}.
	\]
	\(\alpha^2\ne\operatorname{id}\)，所以阶恰为四。在 \(y^2=x^3+1\) 上取本原三次单位根 \(\zeta\)，用 \(u=\zeta\) 得 \(\beta(x,y)=(\zeta^2x,y)\)，阶恰为三。两式都按齐次坐标延拓并固定原点。
\end{solution}

\begin{exercise}\label{b3:ex:H-isogeny-composition}
	设 \(k\) 为代数闭域，\(\operatorname{char}k\ne2,3\)，有同源 \(\varphi:E\rightarrow E'\)、\(\psi:E'\rightarrow E''\)。
	\begin{enumerate}
		\item 证明 \(\psi\circ\varphi\) 的次数为两次数之积，其核的点群有短正合列
		\[
			0\longrightarrow\ker\varphi(k)
			\longrightarrow\ker(\psi\circ\varphi)(k)
			\xrightarrow{\ \varphi\ }\ker\psi(k)\longrightarrow0.
		\]
		\item 证明复合同源在原点的分歧指数为两个原点分歧指数之积，并据此判断复合同源何时可分。求 \([2^r]\) 的次数和核的点数，其中 \(r\ge1\)。
		\item 在 \(k=\overline{\mathbb F}_5\)、\(E:y^2=x^3-x\) 上，令 \(F(x,y)=(x^5,y^5)\)。求 \(F\circ[2]\) 的次数、原点分歧指数和核的点数。
	\end{enumerate}
\end{exercise}
\begin{solution}
	1. 复合保持原点和群运算，且有限满射的复合仍为有限满射，故是同源。函数域扩张的塔式公式给出
	\[
		\deg(\psi\circ\varphi)=\deg\psi\cdot\deg\varphi.
	\]
	题中限制映射的核恰为 \(\ker\varphi(k)\)。对每个 \(Q\in\ker\psi(k)\)，\(\varphi\) 的满射性及 \(k\) 的代数闭性给出 \(P\in E(k)\) 使 \(\varphi(P)=Q\)，于是 \(P\in\ker(\psi\circ\varphi)(k)\)。这证明限制映射也满射。因此短列正合，且
	\[
		\#\ker(\psi\circ\varphi)(k)
		=\#\ker\varphi(k)\cdot\#\ker\psi(k).
	\]

	2. 取 \(E,E',E''\) 原点处的统一化元 \(t,t',t''\)。若 \(\varphi^*t'\) 的阶为 \(e_\varphi\)，\(\psi^*t''\) 的阶为 \(e_\psi\)，则拉回后 \((\psi\circ\varphi)^*t''\) 的阶为 \(e_\varphi e_\psi\)，因为单位在局部同态下仍为单位。故复合的分歧指数为二者之积。由同源的可分性判据，复合可分当且仅当两个指数都为一，也就是两个同源均可分。

	倍点同源的次数为四，分歧指数为一，且有四个核点。\([2^r]\) 是 \([2]\) 的 \(r\) 次复合，因此次数和核的点数均为 \(4^r\)，而分歧指数仍为一。这里只计算基数，不据此直接断定核群在任意整数次数下的完整结构。

	3. 正文已经算出 \(\deg F=5\)、\(e_F=5\)、\(\#\ker F(k)=1\)，以及 \(\deg[2]=4\)、\(e_{[2]}=1\)、\(\#E[2](k)=4\)。所以
	\[
		\deg(F\circ[2])=20,\qquad
		e_{F\circ[2]}=5,\qquad
		\#\ker(F\circ[2])(k)=4.
	\]
	事实上 \(F\) 在 \(k\) 点上单射，因此这个核的点集就是 \(E[2](k)\)；它的四个点在核概形中各有长度五。
\end{solution}

% BEGIN APPENDIX OUTLINE SYNC
% * 状态：初稿已完成；新增范围尚未正式审核。
% END APPENDIX OUTLINE SYNC
